% !Tex root = ../main.tex

\chapter{Motion of a single plasma particle}

The previous chapter was devoted to exploring some fundamental collective behaviours of plasmas, specifically from a statistical, macroscopic point of view. In this chapter, we will develop our understanding of the microscopic behaviour of individual plasma particles when subjected to some specific configuration of external force fields. However, it should be noted that charged particles, for their part, make up currents that modify these electromagnetic fields. In this way, any individual plasma particle interacts not only with the external fields, but also with those \textit{internal} force fields produced by the plasma itself, which consequently makes an influence on the individual particle's dynamics. Since any plasma particle in the grand ensemble can be considered as an \textit{individual particle}, its dynamics is, to a certain degree, influenced by the global internal configuration of the plasma. Therefore, a significant part of plasma's dynamics is governed by its internal configuration at a given moment in time. In our considerations, we will somewhat ignore (more like \textit{oversimplify}) the problem of continuous self-actualisation of the plasma configuration by treating all force fields, external and internal, as \textit{known} and prescribed.

\quad Moreover, our interest in single-particle dynamics is motivated by the fact that many microscopic behaviours of particular plasma species are \textit{lost} in the statistical process of averaging. We will frequently encounter this feature of how statistically averaging over distribution functions of species, as well as over all plasma species themselves, simply \textit{forgets} certain phenomena that only specific subgroups of particles experience, that are triggered by e.g. particles moving at specific resonant velocities or undergoing cyclotron resonance. These resonant phenomena are of great interest because plasmas do not possess a mechanism that can efficiently extinguish such particularities due to their weak collisionality; hence, these phenomena persist and largely influence the macroscopic scale. 

\section{Hamilton-Lagrange formalism applied to a plasma particle}

From a mechanical point of view, the dynamics of particles can be described from two analogous perspectives: Newtonian mechanics, which provides the equations of motion via Newton's laws by relating the particle's instantaneous acceleration to the force field configuration at a given instant; and Hamiltonian mechanics, which provides the equations of motion by considering the Hamiltonian (and Lagrangian) function whose existence we postulate and presume to contain all information of the system's dynamics. When it comes to the system's dynamics, the two perspectives are equivalent; each method gives physically identical equations of motion. However, Newtonian mechanics is preferable when considering the system's dynamics for a short range of time, while Hamiltonian mechanics enables us to describe the system via symmetries and (adiabatically) conserved quantities valid in the long range.

\begin{multicols}{2}

\begin{center}
    (Relativistic) Newtonian mechanics:
    \begin{equation}
        \mathbf{p}=\left(\gamma_{\mathbf{v}}\right)m \mathbf{v}
    \end{equation}
    
    \begin{equation}
        \frac{\mathrm{d}\mathbf{p}}{\mathrm{d}t}=q \left(\mathbf{E}+\mathbf{v}\times \mathbf{B}\right)+\mathbf{F}_{\mathrm{ext}}
    \end{equation}
\end{center}

\columnbreak

\begin{center}
    Hamiltonian-Lagrangian formalism:
    \begin{equation}
        \mathcal{H}=\sum_{i}\mathbf{P}_i\cdot\dot{\mathbf{Q}_i}-\mathcal{L}
    \end{equation}
    \begin{equation}
        \frac{\partial \mathbf{Q}_i}{\partial t}=\frac{\partial \mathcal{H}}{\partial \mathbf{P}_i}\quad\text{and}\quad \frac{\partial \mathbf{P}_i}{\partial t}=-\frac{\partial \mathcal{H}}{\partial \mathbf{Q}_i}
    \end{equation}
\end{center}

\end{multicols}

where we allow ourselves to treat the classical mechanical problem without invoking relativistic effects.

\vspace*{2.5pt}

\quad For simplicity's sake, we shall ignore generalised coordinates for the time being since we treat the general case without the assumption of any symmetries. The Lagrangian for a charged particle in an electromagnetic field reads:

\begin{equation}
    \mathcal{L}(\mathbf{x},\dot{\mathbf{x}},t)=\frac{1}{2}m\dot{\mathbf{x}}^2+q\dot{\mathbf{x}}\cdot \mathbf{A}\left(\mathbf{x},t\right)-q\phi \left(\mathbf{x},t\right)
\end{equation}

The generalised momentum is thus given by the following, and yields the equation of motion when combined with the Euler-Lagrange equation:

\begin{equation*}
    \mathbf{P}=\frac{\partial\mathcal{L}}{\partial \dot{\mathbf{x}}}=m\dot{\mathbf{x}}+q \mathbf{A}\left(\mathbf{x},t\right)\implies \frac{\mathrm{d}\mathbf{P}}{\mathrm{d}t}=\frac{\partial \mathcal{L}}{\partial \mathbf{x}}=q\nabla\left(\dot{\mathbf{x}}\cdot \mathbf{A}(\mathbf{x},t)\right)-q\nabla\phi \left(\mathbf{x},t\right)
\end{equation*}

\begin{equation*}
    \frac{\mathrm{d}\mathbf{P}}{\mathrm{d}t}=m\ddot{\mathbf{x}}+q \frac{\partial \mathbf{A}}{\partial t}+q\left(\dot{\mathbf{x}}\cdot\nabla\right)\mathbf{A}(\mathbf{x},t)
\end{equation*}

\begin{equation*}
    \nabla \left(\dot{\mathbf{x}}\cdot \mathbf{A}\right)=\left(\dot{\mathbf{x}}\cdot\nabla\right)\mathbf{A}+\left(\mathbf{A}\cdot\nabla\right)\dot{\mathbf{x}}+\dot{\mathbf{x}}\times \left(\nabla\times \mathbf{A}\right)+\mathbf{A}\times \left(\nabla\times \dot{\mathbf{x}}\right)=\left(\dot{\mathbf{x}}\cdot\nabla\right)\mathbf{A}+\dot{\mathbf{x}}\times \mathbf{B}
\end{equation*}

where, in the last identity, we used the fact that the velocity of the particle is dependent, and moreover parametrised, by time $t$ in virtue of parametrising its trajectory with time. Combining everything gives:

\begin{equation}
    m\ddot{\mathbf{x}}=q \left(-\nabla\phi(\mathbf{x},t)-\frac{\partial \mathbf{A}(\mathbf{x},t)}{\partial t}\right)+q\dot{\mathbf{x}}\times \mathbf{B}(\mathbf{x},t)
\end{equation}

which is identical to Newton's second law given $\mathbf{E}=-\nabla \phi-\partial \mathbf{A}/\partial t$. The associated Hamiltonian reads:

\begin{equation}
    \mathcal{H}(\mathbf{x},\mathbf{P},t)=\frac{1}{2}m\dot{\mathbf{x}}^2+q\phi(\mathbf{x},t)=\frac{\left(\mathbf{P}-q \mathbf{A}(\mathbf{x},t)\right)^2}{2m}+q\phi(\mathbf{x},t)
\end{equation}

Here, it is important to note that the Lagrangian formalism, specifically the Euler-Lagrange equation itself, is frame- and coordinate-invariant, i.e. the form of the equation is independent of the coordinate system or reference frame which we use to describe our physical system. The Hamiltonian formalism, on the contrary, is indeed coordinate- and frame-dependent. Notably, the Hamiltonian function often corresponds to the energy of the physical system in consideration. Nonetheless, the energy and (generalised) momentum are frame-dependent, and so are the possible associated conservation laws of energy and momentum. Specifically, it should be noted that the Lagrange and the Hamilton formalisms are not entirely the same. In practice, Lagrangian formalism allows us to identify symmetries, and by extension conservation laws, via the appearence of \textit{cyclic} coordinates. Meanwhile, Hamiltonian formalism is heavily based on the concept of canonical transformations which preserve Hamilton's equations of motion, but make it possibly easier to solve these equations by specifically identifying constants of motion. In that way, possessing a constant of motion allows us to \textit{transform away} a specific coordinate and reduce the complexity of the system of equations.

\subsection{Concepts from classical mechanics}

Every physical system possesses at least one constant of motion: initial conditions. Unfortunately, initial conditions as a constant of motion is generally of little usefulness when trying to concretely solve for the system's evolution. Morever, we are not fully guaranteed to be able to find other constants of motion without proper inspection into the equations of motion. The situation is, however, different if, \textit{a priori}, we know that the system's motion is periodic or \textit{almost} periodic. Such periodic (resp. almost periodic) physical systems form closed (resp. almost closed) paths in the phase space $(\mathbf{Q}, \mathbf{P})$. One can then discuss about Poincaré (resp. adiabatic) invariants of the system. Concretely, we can introduce a general adiabatic invariant as the integral of the system's action taken over a single period of motion.

\subsubsection{Action-angle variables and adiabatic invariants}

Let $\mathcal{H}=\mathcal{H}\left(q,p,\lambda\right)=h$ be the analytical Hamiltonian function of a one-degree of freedom physical system with value $h$, where $q$ and $p$ are conjugate variables, and $\lambda$ is the system's hyperparameter which we will later use to mimic the time-evolution of the Hamiltonian itself, i.e. contains information of the explicit time-dependence of the Hamiltonian. Let the variable pair $(q,p)$ belong to an open domain $D$ of a two-dimensional manifold, and $\lambda\in I\subset \mathbb{R}$. Let us assume that the trajectories given by Hamilton's equations are periodic of period $T(h,\lambda)$. We shall at first assume $\lambda$ as time-independent, then proceed with the general assumption of $\lambda=\lambda(t)$

\vspace*{2.5pt}

\quad The canonical transformation from $(q,p)$ to the action-angle variables $(\psi, J)$ is usually defined by means of a generating function $S(q,J,\lambda)$ such that:

\begin{equation*}
    p=\frac{\partial S(q,J,\lambda)}{\partial q}\quad\text{and}\quad \psi=\frac{\partial S(q,J,\lambda)}{\partial J}
\end{equation*}

Notably, this means that $p=p(q,J,\lambda)$ and $\psi=\psi(q,J,\lambda)$. Starting from the principle of least action, in terms of canonical variables $(q,p)$ and $(\psi, J)$, we can write:

\begin{equation*}
    \delta \mathcal{S}=\delta\int_{t_0}^{t}\mathcal{L}\left(q,\dot{q},\lambda\right)\,\mathrm{d}t=\delta\int_{t_0}^{t}\left(p\dot{q}-\mathcal{H}\left(q,p,\lambda\right)\right)\,\mathrm{d}t=0
\end{equation*}

\begin{equation*}
    \delta\mathcal{S}=\delta\int_{t_0}^t\mathcal{L}'(\psi,\dot{\psi},\lambda)\,\mathrm{d}t=\delta\int_{t_0}^{t}\left(J\dot{\psi}-\mathcal{K}\left(\psi,J,\lambda\right)\right)\,\mathrm{d}t=0
\end{equation*}

where $\mathcal{S}$ is the action which we minimise in virtue of the principle of least action. The covariance of the Lagrangian implies that $\mathcal{L}$ and $\mathcal{L}'$ can be related by the total time derivative of a function $F$ (in particular, this function is related to the generating function $S$), such that $\mathcal{L}'=\mathcal{L}+\mathrm{d}S/\mathrm{d}t$. The integrands of the integrals which we ought to minimise read:

\begin{equation}
    \label{H-8}
    p\dot{q}-\mathcal{H}\left(q,p,\lambda\right)=J\dot{\psi}-\mathcal{K}\left(\psi,J,\lambda\right)+\frac{\mathrm{d}F}{\mathrm{d}t}
\end{equation}

Let us impose $F=S(q,J,\lambda)-\psi J$, thus the total time derivative is given by:

\begin{equation*}
    \frac{\mathrm{d}F}{\mathrm{d}t}=\frac{\partial S(q,J,\lambda)}{\partial q}\dot{q}+\frac{\partial S(q,J,\lambda)}{\partial J}\dot{J}-\dot{\psi}J-\psi\dot{J}+\frac{\partial S(q,J,\lambda)}{\partial t}
\end{equation*}

We can now input the time derivative into (\ref{H-8}), which gives, after rearranging:

\begin{equation}
    \label{H-9}
    \left(p-\frac{\partial S(q,J,\lambda)}{\partial q}\right)\dot{q}+\left(\psi-\frac{\partial S(q,J,\lambda)}{\partial J}\right)\dot{J}=\mathcal{H}(q,p,\lambda)-\mathcal{K}(\psi,J,\lambda)+\frac{\partial S(q,J,\lambda)}{\partial t} 
\end{equation}

It is now important to note that $q$ and $J$, and by extension $\dot{q}$ and $\dot{J}$, are independent variables. Therefore, for them to be independent, the terms in parenthesis multiplying each must vanish, so that we indeed recover the aforementioned coordinate transformation:

\begin{equation*}
    p=\frac{\partial S(q,J,\lambda)}{\partial q}\quad\text{and}\quad \psi=\frac{\partial S(q,J,\lambda)}{\partial J}
\end{equation*}

That being said, the right-hand side of (\ref{H-9}) must equally vanish, giving:

\begin{equation}
    \label{H-10}
    \mathcal{K}(\psi,J,\lambda)=\mathcal{H}(q,p,\lambda)+\frac{\partial S(q,J,\lambda)}{\partial t}
\end{equation}

So far, the variables $(\psi, J)$ were generic canonical variables; now, we introduce them as action-angle variables such that, for a time-independent Hamiltonian, we impose that $\mathcal{K}(\psi, J,\lambda)=\mathcal{K}(J,\lambda)$. Indeed, a time-independent Hamiltonian is invariant and describes the energy of the system. It is precisely through the hyperparameter $\lambda$ that we ensure its time-independency. Moreover, if the Hamiltonian is invariant, could we impose that it is identically zero? From (\ref{H-10}), this would mean:

\begin{equation}
    \label{H-11}
    \mathcal{H}\left(q,\frac{\partial S(q,J,\lambda)}{\partial q},\lambda\right)+\frac{\partial S(q,J,\lambda)}{\partial t}=0
\end{equation}

The above equation (\ref{H-11}) is commonly referred to as the Hamilton-Jacobi equation. In general, the Hamilton-Jacobi equation is an implicit partial differential equation that includes only the first-order partial derivatives of the generating function with respect to the canonical and time-like variables; specifically, for an $n$-degree of freedom system, it includes $n+1$ partial derivatives, and thus for a unique solution, we require $n+1$ initial conditions, i.e. constants. One of these constants is surely an additive constant, since $S+\text{const.}$ is also a solution to (\ref{H-11}). We can thus search for a solution of the form $S(q,J,\lambda)=S(q;\alpha;\lambda)$, where $\alpha$ is a constant. Once $S(q;\alpha;\lambda)$ is found, we can retrieve the generating function by imposing $S(q,J,\lambda)=S(q;\alpha=J;\lambda)$.

\vspace*{2.5pt}

\quad Since $\mathcal{H}$ is time-independent, we can search for a solution for $S$ under the form:

\begin{equation}
    \label{H-12}
    S(q;\alpha;\lambda)=W(q;\alpha)-\alpha\lambda\implies\mathcal{H}\left(q,\frac{\partial W(q;\alpha)}{\partial q},\lambda\right)=\alpha\equiv h
\end{equation}

which is commonly called the \textit{characteristic} Hamilton-Jacobi equation. From here, we can try to directly determine $S(q,J,\lambda)$ for a known Hamiltonian, but it is worthwhile to use this relation in order to construct a canonical transformation via the generating function $S(q,J,\lambda)=W(q;h=J)$. Note that since $\alpha$ is equivalent to the value of the Hamiltonian, we will write $\alpha=h$ from now on. Since $W(q;h)$ is time-independent, we have $p=\partial W/\partial q$, $\psi=\partial W/\partial J$, and $K(\psi, J,\lambda)=J$. Since $\psi$ is a cyclic coordinate, we have $J=\text{const.}$ and $\dot{\psi}=\text{const.}$ It is sometimes more advantageous to consider a more general transformation, $S(q,J,\lambda)=W(q;h=\beta(J))$ where the function $\beta$ defines a change of variables:

\begin{equation*}
    p=\frac{\partial W(q;\beta(J)}{\partial q},\quad \psi=\frac{\partial W(q;\beta(J)}{\partial \beta}\frac{\partial \beta}{\partial J},\quad\mathcal{K}(\psi, J)=\beta(J)\quad\implies\quad\dot{\psi}=\frac{\partial\mathcal{K}(\psi,J)}{\partial J}=\frac{\mathrm{d}\beta}{\mathrm{d}J}
\end{equation*}

Suppose that we have now solved for $W(q;h)$; the problem is now reduced to finding $\beta(J)$ such that the generating function ties to a canonical transformation where $J$ is constant and $\psi$ increases by $2\pi$ during a period of motion. We will therefore express the angle-like variable $\psi=\psi(q)$ such that $\psi(t+T(h,\lambda))=\psi(t)$; we can now write the periodic condition on $\psi$ as:

\begin{equation*}
    \Delta_\text{cyc}\psi=\oint_\text{cyc}\frac{\mathrm{d}\psi}{\mathrm{d}q}\,\mathrm{d}q=2\pi\quad\text{with}\quad\psi=\frac{\partial S(q,J,\lambda)}{\partial J}
\end{equation*}

We can now evaluate the total derivate as follows:

\begin{equation*}
    \frac{\mathrm{d}\psi}{\mathrm{d}q}=\frac{\partial^2 S(q,J,\lambda)}{\partial q\partial J}+\frac{\partial^2 S(q,J,\lambda)}{\partial J^2}\frac{\partial J}{\partial q}
\end{equation*}

However, $\partial J/\partial q=0$ since $(q,J)$ are independent variables, the second term vanishes, giving:

\begin{equation*}
    2\pi=\oint_\text{cyc}\frac{\partial^2 S(q,J,\lambda)}{\partial q\partial J}\,\mathrm{d}q=\frac{\partial}{\partial J}\oint_\text{cyc}\frac{\partial S(q,J,\lambda)}{\partial q}\,\mathrm{d}q=\frac{\partial }{\partial J}\oint_\text{cyc}\frac{\partial W(q;h)}{\partial q}\,\mathrm{d}q
\end{equation*}

Since $(q,J)$ are independent, and $h$ a constant, the coordinate $J$ is expressed in terms of $h$ as follows:

\begin{equation}
    J=\frac{1}{2\pi}\oint_\text{cyc}\frac{\partial W(q;h)}{\partial q}\,\mathrm{d}q\implies \beta^{-1}(h)=\frac{1}{2\pi}\oint_\text{cyc}\frac{\partial W(q;h)}{\partial q}\,\mathrm{d}q
\end{equation}

Inverting the above provides us with the function $\beta(J)$ which solves the initial problem.

\vspace*{2.5pt}

\quad The terminology of action-angle variables is now evident; we have that $\dot{\psi}=\mathrm{d}\beta/\mathrm{d}J$, which is independent of time, so that $\psi(t)=\psi(t=0)+\beta'(J)t$, which directly gives access to an expression of the period of motion. The variable $J$ can, on the contrary, be written using the original momentum variable $p$ to give $\frac{1}{2\pi}\oint_\text{cyc}p\,\mathrm{d}q$ via the generating function, which has the same units as action, but it is related to action itself via $\mathcal{S}=J-ht$.

\vspace*{2.5pt}

\quad Moreover, since $J$ is constant along a periodic trajectory related to the total energy $h$ itself, we can interpret $J$ as a label for trajectories. First, periodic orbits are closed curves in the phase space $(q,p)$. Second, since $\mathcal{H}$ is a single-valued function, trajectories with different values for $J$ cannot intersect or touch one another; otherwise, there would exist another trajectory, part of which corresponds to the first, and another part to the second. This being, by assumption, a valid trajectory, it would have to be described by a single value for the Hamiltonian, which provides us with the necessary contradiction. Therefore, curves of constant $J$ form a nested family of trajectories.

\vspace*{2.5pt}

\quad On the other hand, the variable $\psi$ serves as an angle-like variable that allows us to complete the nested family of $J=\text{const}.$ curves into a coordinate system. Since $J$ is can be expressed as a function with the single argument $h$, we can travel \textit{across} trajectories only through curves $\psi=\text{const.}$ This will be important later.

\vspace*{2.5pt}

\quad We are now prepared to determine the generating function $S(q,J,\lambda)$. From the Hamilton-Jacobi equation, we can define the function $P(q,h,\lambda)$ as:

\begin{equation*}
    \mathcal{H}(q,P,\lambda)=h\implies\frac{\partial S(q,J,\lambda)}{\partial q}=P(q,h,\lambda)\implies S(q,J,\lambda)=\int_{q_0(h,\lambda)}^q P(q',h,\lambda)\,\mathrm{d}q'+F(h)
\end{equation*}

where the first term is directly obtained from (\ref{H-12}) when combined with the newly defined function, while $F(h)$ is yet to be determined. This function is designed to force $\psi=0$ on the so-called \textit{curve of initial conditions}, which is parametrised by $h$ as $(q_0(h,\lambda),p_0(h,\lambda))$. Mathematically, we impose the following boundary condition:

\begin{equation}
    \label{H-14}
    0\stackrel{!}{=}\psi(q_0(h,\lambda),J,\lambda)=\frac{\partial S(q,J,\lambda)}{\partial J}\biggr|_{q_0(h,\lambda)}=\frac{\partial S(q,h(J),\lambda)}{\partial h}\biggr|_{q_0(h,\lambda)}\frac{\mathrm{d}h}{\mathrm{d}J}
\end{equation}

Since the frequency $\mathrm{d}h/\mathrm{d}J$ is non-zero, we have that the partial derivative $\partial S/\partial h(q_0(h,\lambda),J,\lambda)$ must evaluate to zero.
We can now compute this partial derivative:

\begin{equation*}
    \frac{\partial S(q,J,\lambda)}{\partial h}=\frac{\partial }{\partial h}\int_{q_0(h,\lambda)}^q P(q',h,\lambda)\,\mathrm{d}q'+\frac{\mathrm{d}F}{\mathrm{d}h}
\end{equation*}

By the Leibniz rule of differentiation under the integral sign, we have:

\begin{equation*}
    \frac{\partial S(q,J,\lambda)}{\partial h}=\int_{q_0(h,\lambda)}^q\frac{\partial P(q',h,\lambda)}{\partial h}\,\mathrm{d}q'-P(q_0(h,\lambda),h,\lambda)\frac{\partial q_0(h,\lambda)}{\partial h}+\frac{\mathrm{d}F}{\mathrm{d}h}
\end{equation*}

Evaluating the above at $q=q_0(h,\lambda)$, we can see that the integral term vanishes, while the remaining terms, along with (\ref{H-14}), give:

\begin{equation}
    \label{H-15}
    \frac{\mathrm{d}F}{\mathrm{d}h}=P(q_0(h,\lambda),h,\lambda)\frac{\partial q_0(h,\lambda)}{\partial h}
\end{equation}

From the definition of the function $P(q,h,\lambda)$ as $\mathcal{H}(q,P,\lambda)=h$, evaluating the above at $q=q_0(h,\lambda)$ gives $P(q_0(h,\lambda),h,\lambda)=p_0(h,\lambda)$. Plugging this into (\ref{H-15}) gives an ordinary differential equation for $F(h)$, which can be solved by integrating:

\begin{equation*}
    \frac{\mathrm{d}F}{\mathrm{d}h}=p_0(h,\lambda)\frac{\partial q_0(h,\lambda)}{\partial h}\implies F(h)=\int_{h^\star}^h p_0(h',\lambda)\frac{\partial q_0(h',\lambda)}{\partial h'}\,\mathrm{d}h'
\end{equation*}

This gives the final expression for the generating function $S(q,J,\lambda)$ as follows:

\begin{equation}
    \label{H-18}
    S(q,J,\lambda)=\int_{q_0(h,\lambda)}^q P(q',h,\lambda)\,\mathrm{d}q'+\int_{h^\star}^h p_0(h',\lambda)\frac{\partial q_0(h',\lambda)}{\partial h'}\,\mathrm{d}h'
\end{equation}

where we remind $h=\beta(J)$.

\vspace*{2.5pt}

\quad Let us now elaborate more on what is meant by concepts of \textit{curve of initial conditions} and similar. The idea of the curve of initial conditions is that, instead of picking an arbitrary point $(q_0,p_0)$ for each trajectory $h=\text{const.}$ which would belong to said trajectory, we can instead use the fact that curves of constant $\psi$ run across trajectories. This allows us to control all possible initial conditions $(q_0,p_0)$ by collecting them within this $\psi=\text{const.}$ curve. Since the choice is arbitrary, we can decide that they all lie on $\psi=0$ curve. Since this curve runs across trajectories, each point $(q_0,p_0)\in\{\psi=0\}$ belongs to a distinct trajectory, thus we can say $(q_0,p_0)=\left(q_0(h,\lambda),p_0(h,\lambda)\right)$, where we include the hyperparameter-dependence because the curve of initial conditions may also depend on time.

\vspace*{2.5pt}

\quad We will assume that the curve of initial conditions is transversal to all periodic trajectories. We say that two intersecting curves are \textit{transveral} if they cross each other at a non-zero angle, i.e. they intersect each other in a way that the tangent vectors along the two curves are never parallel, i.e. never linearly dependent. This condition may be somewhat restrictive but ensures that the curve $\psi=0$ is well-defined, and moreover, that the canonical transformation $(q,p)\mapsto (\psi,J)$ is a valid transformation whose Jacobian is everywhere non-zero. In other words, transversality ensures that the canonical transformation is smooth and inertible everywhere. However, if we are not interested in the effect of $\lambda$ on initial conditions, we can without issue drop the notion of fixing $\psi=0$ to describe the curve of initial conditions, and therefore also drop the second term in (\ref{H-18}).

\vspace*{2.5pt}

\quad Let us now consider that $\lambda$ is no longer a constant but a function of time, $\lambda=\varepsilon t$. The Hamiltonian in $(\psi, J)$ canonical coordinates reads directly from (\ref{H-10}):

\begin{equation}
    \label{H-17}
    \mathcal{H}'(\psi,J,\varepsilon t)=\mathcal{H}(q,p,\varepsilon t)+\frac{\partial S(q,J,\varepsilon t)}{\partial t}=\mathcal{K}(J,\varepsilon t)+\frac{\partial S(q,J,\lambda)}{\partial \lambda}\underbrace{\frac{\mathrm{d}\lambda}{\mathrm{d}t}}_{\varepsilon}
\end{equation}

We can therefore define the last term as the remainder function, namely:

\begin{equation}
    \varepsilon R(\psi, J, \varepsilon t)=\left(\frac{\partial S(q,J,\varepsilon t)}{\partial t}\right)_{q(\psi,J),J}
\end{equation}

where the subscripts denote that the partial derivative is taken such that $q$ and $J$ are kept constant, and then $q=q(\psi, J)$ is expressed in terms of $\psi$ and $J$. The remainder function can be found explicitly through the Hamiltonian; we trivially have that $\frac{\partial \mathcal{H}}{\partial \lambda}=\frac{1}{\varepsilon}\frac{\partial \mathcal{H}}{\partial t}=\frac{1}{\varepsilon}\frac{\mathrm{d}\mathcal{H}}{\mathrm{d}t}$, where the last equatily follows from the fact that the Hamiltonian is in involution with itself. Now, since the value of the old Hamiltonian $\mathcal{H}$ is equal to $h=\mathcal{K}(J,\lambda)$, we obtain:

\begin{equation*}
    \frac{\mathrm{d}\mathcal{H}}{\mathrm{d}t}=\frac{\mathrm{d}\mathcal{K}}{\mathrm{d}t}=\frac{\partial \mathcal{K}}{\partial J}\dot{J}+\frac{\partial\mathcal{K}}{\partial \lambda}\dot{\lambda}=\frac{\partial \mathcal{K}}{\partial J}\dot{J}+\varepsilon \frac{\partial \mathcal{K}}{\partial \lambda}
\end{equation*}

Now, $J$ is no longer guaranteed to be constant. We can find its time evolution equation through $\mathcal{H}'$:

\begin{equation*}
    \dot{J}=-\frac{\partial \mathcal{H}'(\psi, J,\varepsilon t)}{\partial \psi}=-\frac{\partial}{\partial \psi}\left(\mathcal{K}(J,\varepsilon t)+\varepsilon R(\psi, J,\varepsilon t)\right)=-\varepsilon \frac{\partial R(\psi, J,\varepsilon t)}{\partial \psi}
\end{equation*}

Therefore, we have that:

\begin{equation*}
    \frac{\partial \mathcal{H}(q,p,\lambda)}{\partial \lambda}\Bigr|_{\left(q(\psi,J),p(\psi,J)\right)}=-\frac{\partial\mathcal{K}(J,\lambda)}{\partial J}\frac{\partial R(\psi,J,\varepsilon t)}{\partial \psi}+\frac{\partial \mathcal{K}(J,\lambda)}{\partial \lambda}
\end{equation*}

We can now extract the derivative concerning the remainder function:

\begin{equation*}
    \frac{\partial R(\psi, J,\varepsilon t)}{\partial \psi}=\left(\frac{\partial\mathcal{K}(J,\lambda)}{\partial J}\right)^{-1}\left\{\frac{\partial \mathcal{K}(J,\lambda)}{\partial \lambda}-\frac{\partial \mathcal{H}(q(\psi,J),p(\psi, J),\lambda)}{\partial \lambda}\right\}
\end{equation*}

For the sake of clarity, we will drop the functional dependencies. Integrating over $\psi$ gives:

\begin{equation}
    \label{H-19}
    R(\psi, J,\varepsilon t)-R(0,J,\varepsilon t)=\left(\frac{\partial\mathcal{K}}{\partial J}\right)^{-1}\int_0^\psi \left(\frac{\partial \mathcal{K}}{\partial \lambda}-\frac{\partial \mathcal{H}}{\partial \lambda}\right)\,\mathrm{d}\psi'
\end{equation}

where $R(0,J,\varepsilon t)$ denotes the remainder value on the curve of initial conditions, which we can determine from (\ref{H-18}), first by partially differentiating with respect to $\lambda$, and only then inputting $q=q_0(h,\lambda)$. One additional issue that has to be adressed is that $S(q,J,\lambda)$ depends on $h=\mathcal{K}(J,\lambda)$ through $J$, which itself can depend on $\lambda$. When considering $\lambda$ as a constant, this issue was simply inexistent. Luckily, we can quickly resolve the problem:

\begin{equation*}
    R(\psi,J,\lambda)=\left(\frac{\partial S(q,J,\lambda)}{\partial \lambda}\right)_{q,J}=\left(\frac{\partial S(q,J,\lambda)}{\partial \lambda}\right)_{q,h}+\left(\frac{\partial S(q,J,\lambda)}{\partial h}\right)_{q,\lambda}\left(\frac{\partial\mathcal{K}(J,\lambda)}{\partial \lambda}\right)_J
\end{equation*}

where it is understood that $q=q(\psi, J)$ must be expressed in terms of the action-angle variables. It is now useful to remind ourselves that choosing $\psi=0$ ensures that $\partial S/\partial h=0$ at $q=q_0(h,\lambda)$, so that the second term vanishes exactly. We can thus proceed to find the remaining term:

\begin{equation}
    \label{H-20}
    R(\psi,J,\lambda)=\left(\frac{\partial S(q,J,\lambda)}{\partial \lambda}\right)_{q,J}=\int_{q_0}^q \frac{\partial P(q',h,\lambda)}{\partial \lambda}\,\mathrm{d}q'-P(q_0,h,\lambda)\frac{\partial q_0}{\partial \lambda}+\int_{h^\star}^h\left(\frac{\partial p_0}{\partial \lambda}\frac{\partial q_0}{\partial h'}+p_0 \frac{\partial^2 q_0}{\partial h'\partial\lambda}\right)\mathrm{d}h'
\end{equation}

where it is understood that $q=q(\psi,J)$, $q_0=q_0(h,\lambda)$ and $p_0=p_0(h,\lambda)$. The first integral vanishes exactly when evaluated at $q=q_0$, i.e. $\psi=0$. We already know $P(q_0,h,\lambda)=p_0$. The second term in the last integral can be manipulated as follows:

\begin{equation}
    \label{H-21}
    \int_{h^\star}^h p_0(h',\lambda)\frac{\partial^2q_0(h',\lambda)}{\partial h'\partial\lambda}\,\mathrm{d}h'=\left[p_0(h',\lambda)\frac{\partial q_0(h',\lambda)}{\partial \lambda}\right]\biggr|_{h^\star}^h-\int_{h^\star}^h \frac{\partial p_0(h',\lambda)}{\partial h'}\frac{\partial q_0(h',\lambda)}{\partial \lambda}\,\mathrm{d}h'
\end{equation}

By inputting (\ref{H-21}) in (\ref{H-20}), we observe that the two terms $p_0\partial q_0/\partial\lambda(h,\lambda)$ cancel:

\begin{equation}
    R(\psi=0,J,\lambda)=\int_{h^\star}^h \left(\frac{\partial p_0(h',\lambda)}{\partial \lambda}\frac{\partial q_0(h',\lambda)}{\partial h'}-\frac{\partial p_0(h',\lambda)}{\partial h'}\frac{\partial q_0(h',\lambda)}{\partial \lambda}\right)\mathrm{d}h'-p_0(h^\star,\lambda)\frac{\partial q_0(h^\star,\lambda)}{\partial \lambda}
\end{equation}

which is therefore the correction factor $R(0,J,\varepsilon t)$ in (\ref{H-19}) for when the time-dependency of $\lambda$ influences the curve of initial conditions. Note that $R(\psi, J, \lambda)$ is periodic in $\psi$ through the dependence of $S$ on $q$ which, being single-valued, is periodic in $\psi$.

\vspace*{2.5pt}

\quad We can now consider the case when $\lambda$ is a \textit{slow} function of time, such that its derivatives are smaller and smaller, i.e. that there exists $\varepsilon>0$ such that:

\begin{equation*}
    \frac{1}{n!}\left|\frac{\mathrm{d}^n\lambda}{\mathrm{d}t^n}\right|\leq \varepsilon^n\quad\text{(uniformly)}
\end{equation*}

Moreover, we will also assume that the periods of the periodic trajectories for when $\lambda=\text{const.}$ in the domain $D\times I$ of definition of the physical system are uniformly bounded, independently of $\varepsilon$, by some period, $T(h,\lambda)\leq T_0$. We are now prepared to apply a classical perturbation technique, called the Lie transform method, to the Hamiltonian system which we derived in (\ref{H-17}):

\begin{equation*}
    \mathcal{H}'(\psi, J, \varepsilon t) = \mathcal{K}(J,\varepsilon t) + \varepsilon R(\psi, J,\varepsilon t)
\end{equation*}

The idea of the Lie transform method is to use the fact that Hamilton's equations generate the so-called Hamiltonian flow which, for its part, is a canonical transformation. Namely, given initial conditions $(q_0,p_0)$ and a Hamiltonian $\mathcal{H}$, the time-evolution $(q_0,p_0)\to(q(t),p(t))$ obtained by solving Hamilton's equations at time $t$ is a canonical transformation. Using this fact, for a continuous deformation $\lambda(\varepsilon, t)$, we can treat $\varepsilon$ as a time-like coordinate and \textit{evolve} the system by varying $\varepsilon$ from zero to its desired value. To do this, we define an auxiliary Hamiltonian $W(\mathbf{x},\varepsilon)$ (not to confuse with the function $W$ above), where $\mathbf{x}(t,\varepsilon)=(\mathbf{q}(t,\varepsilon),\mathbf{p}(t,\varepsilon))^\top$ is a column vector, and $\mathbf{q}=(q_i)_{i}$, $\mathbf{p}=(p_i)_i$ are phase space vectors. We can then solve Hamilton's equations by considering the perturbation $\varepsilon$ as a time-like coordinate. Let $\mathbf{S}$ be the principal symplectic matrix, so that Hamilton's equations for this new Hamiltonian read:

\begin{equation}
    \label{H-23}
    \mathbf{S}=\begin{pmatrix}
        0_n & 1_n\\ -1_n & 0_n
    \end{pmatrix}\implies\frac{\mathrm{d}\mathbf{x}}{\mathrm{d}\varepsilon}=\mathbf{S}\cdot \frac{\partial W(\mathbf{x},\varepsilon)}{\partial \mathbf{x}}\quad\Longleftrightarrow\quad \frac{\mathrm{d}q_i}{\mathrm{d}\varepsilon}=\frac{\partial W}{\partial p_i}\quad\text{and}\quad \frac{\mathrm{d}p_i}{\mathrm{d}\varepsilon}=-\frac{\partial W}{\partial q_i}
\end{equation}

where $0_n$ and $1_n$ represent the null and unit matrices of order $n$, respectively. Now, we call $X(\mathbf{y}(t),\varepsilon)$ the map which evaluates to $\mathbf{x}(t,\varepsilon)$ of the system of ordinary differential equations (\ref{H-23}) at \textit{time} $\varepsilon$. Then, $\mathbf{y}(t)$ represent the \textit{initial conditions} for $\varepsilon=0$ to attain $\mathbf{x}(t)$ when evolving \textit{along} the perturbation parameter $\varepsilon$. 

\vspace*{2.5pt}

\quad Let us elaborate on what $\mathbf{x}$ and $\mathbf{y}$ really represent; $\mathbf{x}$ represents the true coordinates in the phase space of the physical system which evolve according to the exact Hamiltonian $\mathcal{H}(\psi, J, \varepsilon t)$. Meanwhile, $\mathbf{y}$ is a vector in the same phase space which evolves according to a mathematically simpler Hamiltonian, which we will later determine, and this vector is given via the implicitly defined map $\mathbf{x}(t)=X(\mathbf{y}(t),\varepsilon)$, which is in turn governed by the auxiliary Hamiltonian $W(\mathbf{x},\varepsilon)$. Why would this be mathematically simpler? Well, we will try to construct $W(\mathbf{x},\varepsilon)$ such that the \textit{simpler} Hamiltonian $\overline{\mathcal{H}}(\mathbf{y},\varepsilon)=\mathcal{H}\left(X(\mathbf{y},\varepsilon),\varepsilon\right)$ is less complicated (specifically, independent of the angle $\psi$), so that we can easily find how $\mathbf{y}(t)$ evolves with time. Then, to find the exact coordinates in phase space $\mathbf{x}(t)$ of the system are at some time $t$, we use the map $X$ to transform the determined $\mathbf{y}(t)$ into $\mathbf{x}(t)$.

\vspace*{2.5pt}

\quad In short, we construct the auxiliary Hamiltonian $W(\mathbf{x},\varepsilon)$ which \textit{evoles} along the perturbation parameter, and for a fixed $\mathbf{x}$, we call $\mathbf{y}$ the associated \textit{initial conditions} for the \textit{time-like} evolution of the perturbation parameter from zero to $\varepsilon$. We construct this auxiliary Hamiltonian in such a way that we can determine yet another Hamiltonian $\overline{\mathcal{H}}$, which would govern how $\mathbf{y}$ evolves with time $t$, in such a way that solving for $\mathbf{y}$ is simpler than directly solving for $\mathbf{x}$. Then, for any desired time $t$, we can solve for $\mathbf{y}(t)$, and then determine the physical system's true coordinates in phase space $\mathbf{x}(t)=X(\mathbf{y}(t),\varepsilon)$. Note that the above paragraphs are valid in general phase space parametrisations, which is why we used $(\mathbf{q},\mathbf{p})$ as phase space coordinates. We shall later apply our considerations directly to the action-angle variables, i.e. let $\mathbf{x}=(\psi, J)^\top$ and $\mathbf{y}=(\overline{\psi},\overline{J})^\top$, where $\overline{\psi}$ and $\overline{J}$ are coordinates yet to be fully determined.

\vspace*{2.5pt}

\quad Now, we would like to determine the Hamiltonian $\overline{\mathcal{H}}(\mathbf{y},\varepsilon)$ which governs how $\mathbf{y}$ evolves with time. We assume to possess an analytical Hamiltonian $\mathcal{H}(\mathbf{x},\varepsilon)=\mathcal{H}\left(X(\mathbf{y},\varepsilon),\varepsilon\right)$. Let us first consider a generic function $F(\mathbf{x},\varepsilon)$; its total derivative with respect to $\varepsilon$ along the Hamiltonian flow of $W(\mathbf{x},\varepsilon)$ reads:

\begin{equation*}
    \frac{\mathrm{d}F}{\mathrm{d}\varepsilon}=\frac{\partial F}{\partial \varepsilon}+\frac{\partial F}{\partial \mathbf{q}}\cdot \frac{\mathrm{d}\mathbf{q}}{\mathrm{d}\varepsilon}+\frac{\partial F}{\partial \mathbf{p}}\cdot \frac{\mathrm{d}\mathbf{p}}{\mathrm{d}\varepsilon}=\frac{\partial F}{\partial \varepsilon}+\frac{\partial F}{\partial \mathbf{q}}\cdot \frac{\partial W}{\partial \mathbf{p}}-\frac{\partial F}{\partial \mathbf{p}}\cdot \frac{\partial W}{\partial \mathbf{q}}=\frac{\partial F}{\partial \varepsilon}+\left\{F,W\right\}
\end{equation*}

where $\left\{\cdot,\cdot\right\}$ denotes the Poisson bracket. Therefore, we know that $\frac{\mathrm{d}}{\mathrm{d}\varepsilon}=\frac{\partial }{\partial \varepsilon}+\left\{\cdot,W\right\}$. We can perform the Taylor expansion of the original Hamiltonian with respect to the perturbation parameter:

\begin{equation}
    \mathcal{H}(\mathbf{x},\varepsilon)=\sum_{j\geq0}\frac{\varepsilon^j}{j!}H^0_j(\mathbf{x})
\end{equation}

and, in general, let $\left(H^n_j(\mathbf{x})\right)_{j\geq0}$ be the sequence which generates the complete Taylor expansion of the $n$-th order derivative of the original Hamiltonian with respect to the perturbation parameter:

\begin{equation}
    \frac{\mathrm{d}^n\mathcal{H}(\mathbf{x},\varepsilon)}{\mathrm{d}\varepsilon^n}=\sum_{j\geq0}\frac{\varepsilon^j}{j!}H^n_j(\mathbf{x})
\end{equation}

Meanwhile, since we will construct the auxiliary Hamiltonian $W(\mathbf{x},\varepsilon)$ as an analytical function, let:

\begin{equation*}
    W(\mathbf{x},\varepsilon)=\sum_{\ell\geq0}\frac{\varepsilon^\ell}{\ell!}W_{\ell+1}(\mathbf{x})
\end{equation*}

be the Taylor expansion of $W(\mathbf{x},\varepsilon)$ with respect to $\varepsilon$; the shift in index will be useful later.

\vspace*{2.5pt}

\quad Let us assume that we possess the expansion of $\mathrm{d}^{n-1}H(\mathbf{x},\varepsilon)/\mathrm{d}\varepsilon^{n-1}$, given by:

\begin{equation*}
    \frac{\mathrm{d}^{n-1}\mathcal{H}(\mathbf{x},\varepsilon)}{\mathrm{d}\varepsilon^{n-1}}=\sum_{j\geq0}\frac{\varepsilon^j}{j!}H^{n-1}_j(\mathbf{x})
\end{equation*}

Differentiating with respect to $\varepsilon$ gives:

\begin{equation*}
    \frac{\mathrm{d}^n\mathcal{H}(\mathbf{x},\varepsilon)}{\mathrm{d}\varepsilon^n}=\frac{\partial}{\partial \varepsilon}\left(\frac{\mathrm{d}^{n-1}\mathcal{H}(\mathbf{x},\varepsilon)}{\mathrm{d}\varepsilon^{n-1}}\right)+\left\{\frac{\mathrm{d}^{n-1}\mathcal{H}(\mathbf{x},\varepsilon)}{\mathrm{d}\varepsilon^{n-1}},W(\mathbf{x},\varepsilon)\right\}=\sum_{j\geq0}\frac{\varepsilon^{j}}{j!}H^{n-1}_{j+1}(\mathbf{x})+\left\{\frac{\mathrm{d}^{n-1}\mathcal{H}(\mathbf{x},\varepsilon)}{\mathrm{d}\varepsilon^{n-1}},W(\mathbf{x},\varepsilon)\right\}
\end{equation*}

Using the Taylor expansion of the auxiliary Hamiltonian, we can write:

\begin{equation*}
    \left\{\frac{\mathrm{d}^{n-1}\mathcal{H}(\mathbf{x},\varepsilon)}{\mathrm{d}\varepsilon^{n-1}},W(\mathbf{x},\varepsilon)\right\}=\sum_{k\geq0}\sum_{\ell\geq0}\frac{\varepsilon^k}{k!}\frac{\varepsilon^\ell}{\ell!}\left\{H_k^{n-1}(\mathbf{x}),W_{\ell+1}(\mathbf{x})\right\}=
\end{equation*}

\begin{equation*}
    =\sum_{j\geq0}\varepsilon^j\sum_{\ell=0}^j \frac{\left\{H_{j-\ell}^{n-1}(\mathbf{x}),W_{\ell+1}(\mathbf{x})\right\}}{(j-\ell)!\ell!}=\sum_{j\geq0}\frac{\varepsilon^j}{j!}\sum_{\ell=0}^j\binom{j}{\ell}\left\{H_{j-\ell}^{n-1}(\mathbf{x}),W_{\ell+1}(\mathbf{x})\right\}
\end{equation*}

Recollecting terms, we obtain:

\begin{equation}
    \label{H-26}
    \frac{\mathrm{d}^n\mathcal{H}(\mathbf{x},\varepsilon)}{\mathrm{d}\varepsilon^n}=\sum_{j\geq0}\frac{\varepsilon^j}{j!}\left(H^{n-1}_{j+1}(\mathbf{x})+\sum_{\ell=0}^j\binom{j}{\ell}\left\{H^{n-1}_{j-\ell}(\mathbf{x}),W_{\ell+1}(\mathbf{x})\right\}\right)
\end{equation}

This provides us with a recursion formula for determining series coefficients of higher order derivatives of the original Hamiltonian:

\begin{equation}
    H^n_j(\mathbf{x})=H^{n-1}_{j+1}(\mathbf{x})+\sum_{\ell=0}^j\binom{j}{\ell}\left\{H_{j-\ell}^{n-1}(\mathbf{x}),W_{\ell+1}(\mathbf{x})\right\}
\end{equation}

The Hamiltonian $\overline{\mathcal{H}}(\mathbf{y},\varepsilon)$ can also be expressed as a Taylor series:

\begin{equation}
    \overline{\mathcal{H}}(\mathbf{y},\varepsilon)=\sum_{n\geq0}\frac{\varepsilon^n}{n!}\left[\frac{\mathrm{d}^n\overline{\mathcal{H}}(\mathbf{y},\varepsilon)}{\mathrm{d}\varepsilon^n}\right]_{\varepsilon=0}=\sum_{n\geq0}\frac{\varepsilon^n}{n!}\left[\frac{\mathrm{d}^n\mathcal{H}(\mathbf{x},\varepsilon)}{\mathrm{d}\varepsilon^n}\right]_{\varepsilon=0}
\end{equation}

Since $\mathrm{d}^n\mathcal{H}(\mathbf{x},\varepsilon)/\mathrm{d}\varepsilon^n$ is given by expansion (\ref{H-26}), evaluating at $\varepsilon=0$ leaves only the term $j=0$. Moreover, at $\varepsilon=0$, we know that $\mathbf{x}\big|_{\varepsilon=0}=X(\mathbf{y},\varepsilon=0)=\mathbf{y}$ by definition, therefore:

\begin{equation}
    \label{H-29}
    \overline{\mathcal{H}}(\mathbf{y},\varepsilon)=\sum_{n\geq0}\frac{\varepsilon^n}{n!}{H}^n_0(\mathbf{y})\quad\text{where}\quad H^n_0(\mathbf{y})=H^{n-1}_1(\mathbf{y})+\left\{H^{n-1}_0(\mathbf{y}),W_1(\mathbf{y})\right\}
\end{equation}

In fact, the recursion (\ref{H-29}) can be written as $H^n_0(\mathbf{y})=\Tilde{H}^n_0(\mathbf{y})+\left\{H_0^0(\mathbf{y}),W_n(\mathbf{y})\right\}$, where $\Tilde{H}^n_0$ recollects all lower order terms that would've been previously computed, i.e. depends on $(H_j^0)_{j\leq n}$ and $(W_j)_{j<n}$. This means that at each recursive step, we have two \textit{unknowns}, $H^n_0(\mathbf{y})$ and $W_n(\mathbf{y})$. Since the goal was to construct the generating function $W(\mathbf{x},\varepsilon)$ in such a way that the new Hamiltonian is mathematically simple, we can choose $W_n(\mathbf{y})$ to achieve this goal.

\vspace*{2.5pt}

\quad We can now concretely apply our considerations to $\mathbf{x}=(\psi, J)^\top$, and $\mathbf{y}=(\overline{\psi},\overline{J})^\top$, where the meaning of overlined coordinates will be determined later. Let $[0,2\pi]\times \mathcal{J}$ be the set of admissible $(\psi, J)$ which canonically map to the set $D$. Since we're considering only analytical functions $\mathcal{H}(\psi, J, \varepsilon t)$ of $(\psi, J)$, it is certain that $\mathcal{H}\in\mathcal{C}^\infty([0,2\pi]\times\mathcal{J}\times I)$, so that all derivatives of the Hamiltonian are continuous, and by extension bounded on the closure of the Cartesian product. Moreover, we have already established the fact that the Hamiltonian is $2\pi$-periodic with respect to $\psi$; hence, this Hamiltonian admits an everywhere-convergent Fourier expansion with respect to $\psi$. We can thus extract from this Fourier series the constant function which is independent of $\psi$, and an \textit{oscillatory} function of $\psi$ with zero mean for $\psi\in[0,2\pi)$. More precisely, let $\mathbf{F}$ be the vector space of analytical functions of the phase space coordinates $(\psi, J)$ which are $2\pi$-periodic in $\psi$; let $\mathbf{F}_1$ be the vector space of constant functions with respect to $\psi$, i.e. only dependent on $J$ and $\lambda$, and let $\mathbf{F}_2$ be the vector space of analytical functions of $(\psi, J, \lambda)$ which are $2\pi$-periodic in $\psi$, and furthermore, have zero mean when integrated with respect to $\psi\in[0,2\pi)$. By the properties of the Fourier series, we have that $\mathbf{F}=\mathbf{F}_1\oplus \mathbf{F}_2$, i.e. $\mathbf{F}$ is the direct sum of $\mathbf{F}_1$ and $\mathbf{F}_2$. Hence, we can uniquely decompose $\Tilde{H}^n_0(\psi, J)=\langle\Tilde{H}^n_0\rangle(J)+\{\Tilde{H}^n_0\}(\psi, J)$, so that:

\begin{equation}
    \label{H-30}
    H^n_0(\psi,J,\lambda)=\langle\Tilde{H}^n_0\rangle(J,\lambda)+\{\Tilde{H}^n_0\}(\psi,J,\lambda)+\left\{H^0_0(J,\lambda),W_n(\psi,J,\lambda)\right\}
\end{equation}

We would like to make $H^n_0(\psi,J)=H^n_0(J)$ independent of $\psi$, so that we require the last two terms of the right-hand side (\ref{H-30}) to cancel. Calculating this Poisson bracket is tedious if time-dependence of $\lambda$ is considered, but we can elevate $\lambda$ to the status of a canonical coordinate by defining its conjugate momentum, say $\Lambda$, to effectively take care of the time-dependence by extending the phase space $(\psi, J)$ to $(\psi, J, \lambda, \Lambda)$. The extended Hamiltonian is then $\mathcal{H}''(\psi, J,\lambda, \Lambda)=\mathcal{K}(J,\lambda)+\varepsilon R(\psi, J, \lambda) + \varepsilon \Lambda$. Notice that $H^0_0(J)=\mathcal{K}(J,\lambda)$, i.e. the unperturbed Hamiltonian; hence, the Poisson bracket reads:

\begin{equation*}
    \left\{H^0_0(J,\lambda),W_n(\psi, J,\lambda)\right\}=\left(\frac{\partial \mathcal{K}(J,\lambda)}{\partial \psi}\frac{\partial W_n(\psi, J,\lambda)}{\partial J}-\frac{\partial \mathcal{K}(J,\lambda)}{\partial {J}}\frac{\partial W_n(\psi, J,\lambda)}{\partial \psi}\right)+
\end{equation*}

\begin{equation}
    \label{H-31}
    +\left(\frac{\partial \mathcal{K}(J,\lambda)}{\partial \lambda}\frac{\partial W_n(\psi, J, \lambda)}{\partial \Lambda}-\frac{\partial \mathcal{K}(J,\lambda)}{\partial \Lambda}\frac{\partial W_n(\psi, J,\lambda)}{\partial \lambda}\right)
\end{equation}

Since $\mathcal{K}(J,\lambda)$ is independent of $\psi$ and $\Lambda$, the first and last term of the RHS of (\ref{H-31}) vanish. We can then choose $W_n(\psi, J,\lambda)$ to be independent of $\Lambda$ since we are free to construct it to our desire, so that the only surviving term gives:

\begin{equation}
    \label{H-32}
    \left\{H^0_0(J,\lambda),W_n(\psi,J,\lambda)\right\}=-\frac{\partial \mathcal{K}(J,\lambda)}{\partial J}\frac{\partial W_n(\psi, J,\lambda)}{\partial \psi}
\end{equation}

We notice that $\partial\mathcal{K}(J,\lambda)/\partial J=2\pi / T(h,\lambda)$ simply gives the period of a periodic trajectory, which is assumed bounded. We can again use the fact that $W_n(\psi,J,\lambda)$ is necesarily periodic in $\psi$. Assume that $\left\{H^0_0(J,\lambda),W_n(\psi,J,\lambda)\right\}$ evaluates to a non-zero constant function of $\psi$; then, since $\mathcal{K}(J,\lambda)$ is independent of $\psi$, that would necessarily mean that $W_n(\psi, J,\lambda)$ contains a linear term in $\psi$. Such a linear term is not periodic in $\psi$, which contradicts the premise that $W_n(\psi, J,\lambda)$ is periodic in $\psi$; hence, the image of the operator $\left\{H_0^0(J,\lambda),\cdot\right\}$ includes $\mathbf{F}_2$. Notably, (\ref{H-32}) is called a homological equation, and we shall without proof admit the claim that we can solve for $W_n(\psi, J,\lambda)$ since the RHS belongs to the image of $\left\{H^0_0(J,\lambda),\cdot\right\}$. In this way, we can recursively define functions $W_n(\psi, J,\lambda)$ such that the coefficients $H^n_0(\psi, J,\lambda)=H^n_0(J,\lambda)$ are independent of the angle variable $\psi$.

\vspace*{2.5pt}

\quad We would have to stop this recursion at some finite order $n$ because we cannot be certain that the partial sums $W^{(n)}(\mathbf{y},\varepsilon)=\sum_{j=0}^{n-1}\left(\varepsilon^j/j!\right)W_{j+1}(\mathbf{y})$ converge when $n\to\infty$. We shall now consider $n$ to be the finite order at which we stop this recursion. Moreover, we shall call $(\overline{\psi},\overline{J},\lambda,\overline{\Lambda})$ the \textit{averaged} extended phase space coordinates. Let us call $L(W)$ the operator defined by $L(W):f(\mathbf{y})\mapsto f(X(\mathbf{y},\varepsilon))$. Considering $f$ to be the identity function, we have ${L(W)}{\operatorname{{Id}}{(\mathbf{y})}}=\operatorname{{Id}}{\left(X(\mathbf{y},\varepsilon)\right)}=\mathbf{x}$. Therefore, $(\psi, J)=L(W)\mathrm{Id}(\overline{\psi}, \overline{J})$. In fact, this operator executes the recursive Taylor expansion in terms of $\varepsilon$, that is:

\begin{equation}
    J=L(W)\overline{J}=\sum_{j\geq0}\frac{\varepsilon^j}{j!}\sum_{\ell=0}^j\binom{j}{\ell}\left\{\overline{J},W_{\ell+1}(\overline{\psi},\overline{J})\right\}
\end{equation}

We can therefore say that $J=\overline{J}+\varepsilon J_n(\overline{\psi},\overline{J},\varepsilon)=\overline{J}+\varepsilon \left\{\overline{J},W_1(\overline{\psi},\overline{J},\varepsilon)\right\}+\mathcal{O}(\varepsilon^2)$, and explicitly reads:

\begin{equation}
    \label{H-34}
    J=\overline{J}-\varepsilon \left(\frac{\partial \mathcal{K}(\overline{J},\varepsilon t)}{\partial J}\right)^{-1}\left(R(\overline{\psi},\overline{J},\lambda)-\left\langle R(\overline{\psi},\overline{J},\lambda\right\rangle\right)+\mathcal{O}(\varepsilon^2)
\end{equation}

And then the \textit{averaged} coordinate $\overline{\psi}$ is defined as the conjugate coordinate to $\overline{J}$ in the Hamiltonian:

\begin{equation}
    \label{H-35}
    \overline{\mathcal{H}}(\overline{\psi},\overline{J},\lambda,\overline{\Lambda})=\mathcal{K}(\overline{J},\lambda)+\varepsilon\overline{\Lambda}+\sum_{j=1}^n \frac{\varepsilon^j}{j!}\overline{H}^j_0(\overline{J},\lambda)+\varepsilon^{n+1}\overline{H}_r(\overline{\psi},\overline{J},\lambda,\varepsilon)
\end{equation}

In other words $(\overline{\psi},\overline{J})$ are the \textit{initial conditions} for the time-like evolution of the perturbation parameter, $\mathbf{y}(t)$. Since we cut off the recursion to some finite order $n$, these averaged coordinates $(\overline{\psi},\overline{J})$ will depend on $n$; in truth, we would need to explicitly write $(\overline{\psi}_n,\overline{J}_n)$. The Hamiltonian (\ref{H-35}) gives us the time evolution equation for the averaged action variable:

\begin{equation}
    \frac{\mathrm{d}\overline{J}}{\mathrm{d}t}=-\frac{\partial\overline{\mathcal{H}}(\overline{\psi},\overline{J},\lambda,\overline{\Lambda})}{\partial \overline{\psi}}=-\varepsilon^{n+1}\frac{\partial \overline{H}_r(\overline{\psi},\overline{J},\lambda,\varepsilon)}{\partial \overline{\psi}}
\end{equation}

We can now let $\overline{D}$ to be the image of $D$ by the transformation $(q,p)\mapsto(\psi, J)\mapsto(\overline{\psi},\overline{J})$, so that the domain of definition of the system is $\overline{D}\times I$. Furthermore, $\overline{D}$ can be decomposed as $\overline{D}=[0,2\pi)\times[\overline{J}_\text{min},\overline{J}_\text{max}]$. Then, since we assume everything to be analytical, we can say that $\exists C\geq 0:C=\sup_{\overline{D}\times I}\left(\partial\overline{H}_r/\partial\overline{\psi}\right)$. This gives us the final inequality for the averaged action $\overline{J}$:

\begin{equation}
    \label{H-37}
    \left|\frac{\mathrm{d}\overline{J}}{\mathrm{d}t}\right|=\varepsilon^{n+1}\left|\frac{\partial \overline{H}_r}{\partial \overline{\psi}}\right|\leq C \varepsilon^{n+1}\quad\implies\quad \left|\overline{J}(t)-\overline{J}(0)\right|\leq C \varepsilon^{n+1}|t|
\end{equation}

The inequality holds as long as $\lambda=\varepsilon t\in I$ and $\overline{J}\in[\overline{J}_\text{min},\overline{J}_\text{max}]$. However, by (\ref{H-37}), we can ensure that the latter holds so long as:

\begin{equation}
    \overline{J}_\text{min}+C_0 \varepsilon\leq \overline{J}_0 \leq \overline{J}_\text{max}-C_0 \varepsilon\quad\text{for some }C_0\text{ and}\quad -\varepsilon^{-n}\leq t \leq \varepsilon^{-n}
\end{equation}

It is precisely this condition which will provide us with a precise definition of an adiabatic invariant; we say that $\overline{J}(\overline{\psi},J,\lambda)$ is an adiabatic invariant of order $n$ if:

\begin{equation}
    \left|\overline{J}(t)-\overline{J}(0)\right|\leq C \varepsilon^{n+1}\left|t\right|\quad\text{for}\quad \left|t\right|\leq \varepsilon^{-n}
\end{equation}

Following (\ref{H-34}), we can use the triangle inequality and finally obtain:

\begin{equation*}
    \left|J(t)-J(0)\right|\leq \left|J(t)-\overline{J}(t)\right|+\left|\overline{J}(t)-\overline{J}(0)\right|+\left|\overline{J}(0)-J(0)\right|
\end{equation*}

\begin{equation}
    \left|J(t)-J(0)\right|\leq C'\varepsilon\quad\text{for}\quad \left|t\right|\leq \varepsilon^{-n}
\end{equation}

where indeed $J=\frac{1}{2\pi}\oint p\,\mathrm{d}q$ is the \textit{usual} but ill-defined adiabatic invariant. Notice however that all constants, $C$, $C_0$ and $C'$, depend on the order $n$ and can thus attain very large values for larger $n$.

\subsubsection{Interlude to magnetic confinement and the magnetic moment}

We can apply the Hamiltonian formalism and our considerations on general adiabatic invariants to a charged particle moving in an adiabatically non-uniform magnetic field. Let us consider specifically an axisymmetric time-independent magnetic field; the system in consideration is thus a free charged particle of charge $q$ and mass $m$, and we will refer to the magnetic field configuration as the \textit{confinement device}, since such time-independent non-uniform configurations are in practice maintained through currents external to the system. We can choose a gauge for the magnetic vector potential such that:

\begin{equation}
    \mathbf{A}(x,y,z)=\frac{B_0}{2}\left[1+b(z)\right]\left(x\,\mathbf{\hat{e}}_y-y\,\mathbf{\hat{e}}_x\right)
\end{equation}

where the non-uniformity along the axis of the confinement device is captured in the function $b(z)$. The magnetic field is then given by the curl of the magnetic vector potential, that is:

\begin{equation}
    \mathbf{B}(x,y,z)=B_0 \left[1+b(z)\right]\mathbf{\hat{e}}_z-\frac{B_0}{2}\frac{\mathrm{d}b}{\mathrm{d}z}\left(x\,\mathbf{\hat{e}}_x+y\,\mathbf{\hat{e}}_y\right)
\end{equation}

Indeed, we have $\nabla\cdot{\mathbf{B}}=0$, so that any admissible axisymmetric magnetic field configuration is fully determined through $b(z)$. Since the magnetic field is assumed time-independent, any electric field in the confinement device is a pure electrostatic field. The full Hamiltonian of the particle reads:

\begin{equation}
    \mathcal{H}(x,y,z,p_x,p_y,p_z)=\frac{1}{2m}\left(p_x+\frac{qB_0}{2}\left[1+b(z)\right]y\right)^2+\frac{1}{2m}\left(p_y-\frac{qB_0}{2}\left[1+b(z)\right]x\right)^2+\frac{p_z^2}{2m}
\end{equation}

We can introduce the cyclotron frequency as $\Omega(z)=\left(qB_0/m\right) \left[1+b(z)\right]$ for clarity:

\begin{equation}
    \mathcal{H}(x,y,z,p_x,p_y,p_z)=\frac{1}{2m}\left(p_x+\frac{m\Omega(z)}{2}y\right)^2+\frac{1}{2m}\left(p_y-\frac{m\Omega(z)}{2}x\right)^2+\frac{p_z^2}{2m}
\end{equation}

For the time being, we will assume $b(z)>-1$, since otherwise the magnetic field would need to become completely radial with respect to the device's axis in order to subsequently reverse direction. Such situations will be considered later. It is worth noting that this Hamiltonian system is autonomous, and the perturbative time-dependence from the previous section is entirely analogous to the $z$-coordinate, while the perturbation hyperparameter $\lambda$ appears as the function $b$. The velocities along the $x$- and $y$-coordinates are given by:

\begin{equation}
    mv_x=p_x+\frac{m\Omega(z)}{2}y\quad\text{and}\quad mv_y=p_y-\frac{m\Omega(z)}{2}x
\end{equation}

We know \textit{a priori} that our system, i.e. the particle's motion, is almost periodic in the $x$- and $y$-coordinates, which we will call the \textit{transverse} coordinates from now on, and they will together be denoted with $\perp$ in the subscript/superscript. The particle's motion will therefore be composed of gyrating motion (i.e. gyration), gyrating centre motion (i.e. guiding centre motion) and longitudinal motion along the $z$-axis. In the totally uniform case, the guiding centre follows exactly the magnetic field lines (which are straight lines in that case); we can determine the coordinates of the guiding centre as follows:

\begin{equation}
    X_c=x-\rho_x(z)=x+\frac{v_y}{\Omega(z)}=\frac{x}{2}+\frac{p_y}{m\Omega(z)}\quad\text{and}\quad Y_c=y-\rho_y(z)=y-\frac{v_x}{\Omega(z)}=\frac{y}{2}-\frac{p_x}{m\Omega(z)}
\end{equation}

where $\rho_x=-v_y/\Omega(z)$ and $\rho_y=v_x/\Omega(z)$ are the $x$- and $y$-components of the gyroradius (i.e. Larmor radius). Let us now heuristically evaluate the Poisson bracket $\left\{X_c,Y_c\right\}$:

\begin{equation*}
    \left\{X_c,Y_c\right\}=\left(\frac{\partial X_c}{\partial x}\frac{\partial Y_c}{\partial p_x}-\frac{\partial X_c}{\partial p_x}\frac{\partial Y_c}{\partial x}\right)+\left(\frac{\partial X_c}{\partial y}\frac{\partial Y_c}{\partial p_y}-\frac{\partial X_c}{\partial p_y}\frac{\partial Y_c}{\partial y}\right)+\left(\frac{\partial X_c}{\partial z}\frac{\partial Y_c}{\partial p_z}-\frac{\partial X_c}{\partial p_z}\frac{\partial Y_c}{\partial z}\right)=
\end{equation*}

\begin{equation*}
    =\left(-\frac{1}{2m\Omega(z)}\right)+\left(-\frac{1}{2m\Omega(z)}\right)+0=-\frac{1}{m\Omega(z)}
\end{equation*}

This means that we can absorb $-m\Omega(z)$ into one of the guiding centre coordinates in order to construct a conjugate variable pair. Let us set $Y_c$ as a canonical \textit{position} variable, and $m\Omega(z)X_c$ as its conjugate canonical \textit{momentum} variable, so $\left\{Y_c, m\Omega(z)X_c\right\}=1$. We would now need to find another conjugate variable pair in order to construct a canonical transformation for the transverse phase space coordinates; we find our inspiration in the gyroradius components:

\begin{equation*}
    \left\{\rho_x,\rho_y\right\}=\left(\frac{\partial \rho_x}{\partial x}\frac{\partial \rho_y}{\partial p_x}-\frac{\partial \rho_x}{\partial p_x}\frac{\partial \rho_y}{\partial x}\right)+\left(\frac{\partial \rho_x}{\partial y}\frac{\partial \rho_y}{\partial p_y}-\frac{\partial \rho_x}{\partial p_y}\frac{\partial \rho_y}{\partial y}\right)+\left(\frac{\partial \rho_x}{\partial z}\frac{\partial \rho_y}{\partial p_z}-\frac{\partial \rho_x}{\partial p_z}\frac{\partial \rho_y}{\partial z}\right)=
\end{equation*}

\begin{equation*}
    =\left(\frac{1}{2m\Omega(z)}\right)+\left(\frac{1}{2m\Omega(z)}\right)+0=\frac{1}{m\Omega(z)}
\end{equation*}

Hence, $\sqrt{m\Omega(z)}\rho_x$ can act as a canonical \textit{position} variable, while $\sqrt{m\Omega(z)}\rho_y$ can act as its conjugate \textit{momentum} variable. Note that the $z$-dependence of $\Omega(z)$ does not affect the canonical variable transformation $(x,p_x,y,p_y)\mapsto(Y_c,m\Omega(z)X_c,\sqrt{m\Omega(z)}\rho_x,\sqrt{m\Omega(z)}\rho_y)$ since the possible derivatives of $\Omega(z)$ vanish because none of the transverse quantities depend on $p_z$. We repeat that by assuming $b(z)>-1$, the quantity $\sqrt{m\Omega(z)}$ is real-valued and positive.

\vspace*{2.5pt}

\quad Let us call $Y_c=Q_c$, $m\Omega(z)X_c=P_c$, $\sqrt{m\Omega(z)}\rho_x=Q_g$ and $\sqrt{m\Omega(z)}\rho_y=P_g$ with subscripts indicating either relation to the guiding centre with $c$ or gyration with $g$. We can trace back the substitutions for $v_x$ and $v_y$ to obtain the expression of the transverse Hamiltonian:

\begin{equation*}
    v_x=\Omega(z)\rho_y=\sqrt{\frac{\Omega(z)}{m}}P_g\quad\text{and}\quad v_y=-\Omega(z)\rho_x=-\sqrt{\frac{\Omega(z)}{m}}Q_g
\end{equation*}

\begin{equation}
    \label{H-47}
    \mathcal{H}_\perp(Q_c,P_c,Q_g,P_g)=\frac{\Omega(z)}{2}\left(P_g^2+Q_g^2\right)
\end{equation}

The forward and backward substitutions are given by the following relations:

\begin{align*}
    &P_c=m\Omega(z)\left(\frac{x}{2}+\frac{p_y}{m\Omega(z)}\right)& &x=\frac{P_c}{m\Omega(z)}+\frac{Q_g}{\sqrt{m\Omega(z)}}\\
    &Q_c=\frac{y}{2}-\frac{p_x}{m\Omega(z)}& &\frac{p_x}{m\Omega(z)}=\frac{1}{2\sqrt{m\Omega(z)}}P_g-\frac{1}{2}Q_c\\
    &Q_g=\sqrt{m\Omega(z)}\left(\frac{x}{2}-\frac{p_y}{m\Omega(z)}\right)& &y=Q_c+\frac{P_g}{\sqrt{m\Omega(z)}}\\
    &P_g=\sqrt{m\Omega(z)}\left(\frac{y}{2}+\frac{p_x}{m\Omega(z)}\right) & &\frac{p_y}{m\Omega(z)}=\frac{1}{2m\Omega(z)}P_c-\frac{1}{2\sqrt{m\Omega(z)}}Q_g
\end{align*}

It is now necessary to transform the third canonical conjugate variable pair $(z,p_z)$ since the canonical variables $(Q_c,P_c,Q_g,P_g)$ inherently depend on $z$. Let us set the third position variable simply as $Q_z=z$; the third momentum variable can be written as:

\begin{equation*}
    P_z=p_z+F(x,y,z,p_x,p_y)
\end{equation*}

This new third momentum variable must satisfy the canonical relations, which will give five partial differential equations with respect to $p_z, x, p_x, y, p_y$. The canonical relation $\left\{Q_z,P_z\right\}=1$ is already satisfied since we have separated the variable $p_z$ in the above equation. Note that for the sake of conciseness, we will not evaluate all of these terms explicitly. In the end, we would obtain:

\iffalse

\begin{equation}
    \label{H-48}
    \left\{Q_z,P_z\right\}=\left(\frac{\partial Q_z}{\partial x}\frac{\partial P_z}{\partial p_x}-\frac{\partial Q_z}{\partial p_x}\frac{\partial P_z}{\partial x}\right)+\left(\frac{\partial Q_z}{\partial y}\frac{\partial P_z}{\partial p_y}-\frac{\partial Q_z}{\partial p_y}\frac{\partial P_z}{\partial y}\right)+\left(\frac{\partial Q_z}{\partial z}\frac{\partial P_z}{\partial p_z}-\frac{\partial Q_z}{\partial p_z}\frac{\partial P_z}{\partial z}\right)=1
\end{equation}

\begin{equation}
    \label{H-49}
    \left\{Q_c,P_z\right\}=\left(\frac{\partial Q_c}{\partial x}\frac{\partial P_z}{\partial p_x}-\frac{\partial Q_c}{\partial p_x}\frac{\partial P_z}{\partial x}\right)+\left(\frac{\partial Q_c}{\partial y}\frac{\partial P_z}{\partial p_y}-\frac{\partial Q_c}{\partial p_y}\frac{\partial P_z}{\partial y}\right)+\left(\frac{\partial Q_c}{\partial z}\frac{\partial P_z}{\partial p_z}-\frac{\partial Q_c}{\partial p_z}\frac{\partial P_z}{\partial z}\right)=0
\end{equation}

\begin{equation}
    \label{H-50}
    \left\{P_c,P_z\right\}=\left(\frac{\partial P_c}{\partial x}\frac{\partial P_z}{\partial p_x}-\frac{\partial P_c}{\partial p_x}\frac{\partial P_z}{\partial x}\right)+\left(\frac{\partial P_c}{\partial y}\frac{\partial P_z}{\partial p_y}-\frac{\partial P_c}{\partial p_y}\frac{\partial P_z}{\partial y}\right)+\left(\frac{\partial P_c}{\partial z}\frac{\partial P_z}{\partial p_z}-\frac{\partial P_c}{\partial p_z}\frac{\partial P_z}{\partial z}\right)=0
\end{equation}

\begin{equation}
    \label{H-51}
    \left\{Q_g,P_z\right\}=\left(\frac{\partial Q_g}{\partial x}\frac{\partial P_z}{\partial p_x}-\frac{\partial Q_g}{\partial p_x}\frac{\partial P_z}{\partial x}\right)+\left(\frac{\partial Q_g}{\partial y}\frac{\partial P_z}{\partial p_y}-\frac{\partial Q_g}{\partial p_y}\frac{\partial P_z}{\partial y}\right)+\left(\frac{\partial Q_g}{\partial z}\frac{\partial P_z}{\partial p_z}-\frac{\partial Q_g}{\partial p_z}\frac{\partial P_z}{\partial z}\right)=0
\end{equation}

\begin{equation}
    \label{H-52}
    \left\{P_g,P_z\right\}=\left(\frac{\partial P_g}{\partial x}\frac{\partial P_z}{\partial p_x}-\frac{\partial P_g}{\partial p_x}\frac{\partial P_z}{\partial x}\right)+\left(\frac{\partial P_g}{\partial y}\frac{\partial P_z}{\partial p_y}-\frac{\partial P_g}{\partial p_y}\frac{\partial P_z}{\partial y}\right)+\left(\frac{\partial P_g}{\partial z}\frac{\partial P_z}{\partial p_z}-\frac{\partial P_g}{\partial p_z}\frac{\partial P_z}{\partial z}\right)=0
\end{equation}

\begin{equation*}
    \frac{\partial Q_c}{\partial x}=0,\quad \frac{\partial Q_c}{\partial p_x}=-\frac{1}{m\Omega(z)},\quad \frac{\partial Q_c}{\partial y}=\frac{1}{2},\quad \frac{\partial Q_c}{\partial p_y}=0,\quad \frac{\partial Q_c}{\partial z}=+\frac{p_x}{m\Omega^2(z)}\frac{\mathrm{d}\Omega}{\mathrm{d}z}
\end{equation*}

\begin{equation*}
    \frac{\partial P_c}{\partial x}=\frac{m\Omega(z)}{2},\quad \frac{\partial P_c}{\partial p_x}=0,\quad \frac{\partial P_c}{\partial y}=0,\quad \frac{\partial P_c}{\partial p_y}=1,\quad \frac{\partial P_c}{\partial z}=m \frac{\mathrm{d}\Omega}{\mathrm{d}z}\left(\frac{x}{2}+\frac{p_y}{m\Omega(z)}\right)-m\Omega(z)\frac{p_y}{m\Omega^2(z)}\frac{\mathrm{d}\Omega}{\mathrm{d}z}
\end{equation*}

\begin{equation*}
    \frac{\partial Q_g}{\partial x}=\frac{\sqrt{m\Omega(z)}}{2},\quad \frac{\partial Q_g}{\partial p_x}=0,\quad \frac{\partial Q_g}{\partial y}=0,\quad \frac{\partial Q_g}{\partial p_y}=-\frac{1}{\sqrt{m\Omega(z)}}
\end{equation*}

\begin{equation*}
    \frac{\partial P_g}{\partial x}=0,\quad \frac{\partial P_g}{\partial p_x}=\frac{1}{\sqrt{m\Omega(z)}},\quad \frac{\partial P_g}{\partial y}=\frac{\sqrt{m\Omega(z)}}{2},\quad \frac{\partial P_g}{\partial p_y}=0
\end{equation*}

\begin{equation*}
    \frac{\partial Q_g}{\partial z}=\frac{\partial}{\partial z}\left(-\frac{\sqrt{m\Omega(z)}}{\Omega(z)}\left[\frac{p_y}{m}-\frac{\Omega(z)}{2}x\right]\right)=\frac{\partial}{\partial z}\left(-\frac{p_y}{\sqrt{m\Omega(z)}}+\frac{\sqrt{m\Omega(z)}}{2}x\right)
\end{equation*}

\begin{equation*}
    \frac{\partial P_g}{\partial z}=\frac{\partial}{\partial z}\left(\frac{\sqrt{m\Omega(z)}}{\Omega(z)}\left[\frac{p_x}{m}+\frac{\Omega(z)}{2}y\right]\right)=\frac{\partial}{\partial z}\left(\frac{p_x}{\sqrt{m\Omega(z)}}+\frac{\sqrt{m\Omega(z)}}{2}y\right)
\end{equation*}

\begin{equation}
    \label{H-49-2}
    \frac{1}{m\Omega(z)}\frac{\partial F}{\partial x}+\frac{1}{2}\frac{\partial F}{\partial p_y}-\frac{\partial}{\partial z}\left(\frac{p_x}{m\Omega(z)}\right)=0
\end{equation}

\begin{equation}
    \label{H-50-2}
    \frac{m\Omega(z)}{2}\frac{\partial F}{\partial p_x}-\frac{\partial F}{\partial y}+\frac{\partial}{\partial z}\left(\frac{m\Omega(z)}{2}x\right)=0
\end{equation}

\begin{equation}
    \label{H-51-2}
    \frac{\sqrt{m\Omega(z)}}{2}\frac{\partial F}{\partial p_x}+\frac{1}{\sqrt{m\Omega(z)}}\frac{\partial F}{\partial y}+\frac{\partial}{\partial z}\left(-\frac{p_y}{\sqrt{m\Omega(z)}}+\frac{\sqrt{m\Omega(z)}}{2}x\right)=0
\end{equation}

\begin{equation}
    \label{H-52-2}
    -\frac{1}{\sqrt{m\Omega(z)}}\frac{\partial F}{\partial x}+\frac{\sqrt{m\Omega(z)}}{2}\frac{\partial F}{\partial p_y}+\frac{\partial}{\partial z}\left(\frac{p_x}{\sqrt{m\Omega(z)}}+\frac{\sqrt{m\Omega(z)}}{2}y\right)=0
\end{equation}

\begin{align*}
    &\frac{1}{m\Omega(z)}\frac{\partial F}{\partial x}+\frac{1}{2}\frac{\partial F}{\partial p_y}+\frac{p_x}{m}\frac{\Omega'(z)}{\Omega^2(z)}=0\\
    &\frac{1}{m\Omega(z)}\frac{\partial F}{\partial x}-\frac{1}{2}\frac{\partial F}{\partial p_y}+\frac{p_x}{m}\frac{\Omega'(z)}{2\Omega^{2}(z)}-\frac{y}{4}\frac{\Omega'(z)}{\Omega^2(z)}=0
\end{align*}

\begin{align*}
    &\frac{m\Omega(z)}{2}\frac{\partial F}{\partial p_x}-\frac{\partial F}{\partial y}+\frac{\partial}{\partial z}\left(\frac{m\Omega(z)}{2}x\right)=0\\
    &\frac{m\Omega(z)}{2}\frac{\partial F}{\partial p_x}+\frac{\partial F}{\partial y}+\sqrt{m\Omega(z)}\frac{\partial}{\partial z}\left(-\frac{p_y}{\sqrt{m\Omega(z)}}+\frac{\sqrt{m\Omega(z)}}{2}x\right)=0
\end{align*}

\fi

\begin{align*}
    &\frac{\partial F}{\partial x}+\frac{3p_x}{4}\frac{\Omega'(z)}{\Omega(z)}-\frac{m\Omega'(z)y}{8}=0 & &\frac{\partial F}{\partial p_x}+\frac{3x}{4}\frac{\Omega'(z)}{\Omega(z)}+\frac{p_y}{2m}\frac{\Omega'(z)}{\Omega^2(z)}=0\\
    &\frac{\partial F}{\partial p_y}+\frac{p_x}{2m}\frac{\Omega'(z)}{\Omega^2(z)}+\frac{y}{4}\frac{\Omega'(z)}{\Omega(z)}=0 & &\frac{\partial F}{\partial y}+\frac{p_y}{4}\frac{\Omega'(z)}{\Omega(z)}-\frac{m\Omega'(z)x}{8}=0
\end{align*}

We can evaluate all mixed second-order partial derivatives of $F$ with respect to its arguments and find that the corresponding mixed partial derivatives match, meaning that $F$ has an exact differential. What remains to be done is simply integrating each equation to obtain $F(x,y,z,p_x,p_y)$:

\begin{equation*}
    F(x,y,z,p_x,p_y)=m\left[\frac{1}{8}{xy}-\frac{3}{4}x\frac{p_x}{m\Omega(z)}-\frac{1}{4}y\frac{p_y}{m\Omega(z)}-\frac{1}{2}\frac{p_x}{m\Omega(z)}\frac{p_y}{m\Omega(z)}\right]\frac{\mathrm{d}\Omega}{\mathrm{d}z}
\end{equation*}

We would now like to express this function in terms of the new coordinates:

\begin{equation}
    F(Q_c,P_c,Q_g,P_g,z)=\frac{1}{2\Omega(z)}\left[Q_cP_c+\sqrt{m\Omega(z)}\,Q_cQ_g-\frac{P_cP_g}{\sqrt{m\Omega(z)}}\right]\frac{\mathrm{d}\Omega}{\mathrm{d}z}
\end{equation}

We can finally conclude for the new Hamiltonian using $p_z=P_z-F(Q_c,P_c,Q_g,P_g,z)$:

\begin{equation*}
    \mathcal{H}=\frac{\Omega(z)}{2}\left(P_g^2+Q_g^2\right)+\frac{1}{2m}\left[P_z-\frac{\Omega'(z)}{2\Omega(z)}\left(Q_cP_c+\sqrt{m\Omega(z)}\,Q_cQ_g-\frac{P_cP_g}{\sqrt{m\Omega(z)}}\right)\right]^2
\end{equation*}

where the arguments of the Hamiltonian are expressed in terms of $(Q_c,P_c,Q_g,P_g,z,P_z)$. We can immediately recognise that the first term, the transverse Hamiltonian, is equivalent to the Hamiltonian of a harmonic oscillator. We are thus motivated to set $(Q_g,P_g)=(\sqrt{2G}\sin{g},\sqrt{2G}\cos{g})$ where $G$ is an action-like variable and $g$ an angle-like variable. We will refer to $g$ as the gyroangle:

\begin{equation}
    \label{H-49}
    \mathcal{H}(Q_c,P_c,g,G,z,P_z)=\Omega(z)G+\frac{P_z^2}{2m}-\frac{P_z}{m}\Tilde{F}(Q_c,P_c,g,G,z)+\frac{\Tilde{F}^2(Q_c,P_c,g,G,z)}{2m}
\end{equation}

\begin{equation*}
    \Tilde{F}(Q_c,P_c,g,G,z)=\frac{1}{2\Omega(z)}\left(Q_cP_c+\sqrt{2m\Omega(z)G}Q_c\sin{g}-\sqrt{\frac{2G}{m\Omega(z)}}P_c\cos{g}\right)\frac{\mathrm{d}\Omega}{\mathrm{d}z}
\end{equation*}

Let us now assume that $b(z)$ can be written as $\beta(\varepsilon z)$ such that $\beta$ admits no constant coefficient in its Taylor series (we can always rescale the magnetic field strength $B_0$ if $\left\|\mathbf{B}(z=0)\right\|\neq B_0$). The unperturbed Hamiltonian and the perturbation contribution thus read:

\begin{equation*}
    \mathcal{K}(G,z,P_z)=\frac{qB_0}{m}G+\frac{P_z^2}{2m}
\end{equation*}

\iffalse

\begin{equation*}
    \varepsilon R(\varepsilon;Q_,cP_c,g,G,z,P_z)=\frac{qB_0}{m}G\beta(\varepsilon z)+\frac{\Tilde{F}^2(\varepsilon; Q_c,P_c,g,G,z)}{2m}-\frac{P_z}{m}\Tilde{F}(\varepsilon; Q_c,P_c,g,G,z)
\end{equation*}

\begin{equation*}
    \Tilde{F}(\varepsilon; Q_c,P_c,g,G,z)=\frac{\varepsilon\beta'(\varepsilon z)}{2\left[1+\beta(\varepsilon z)\right]}\left(Q_cP_c+\sqrt{2qB_0G \left[1+\beta(\varepsilon z)\right]}Q_c\sin{g}-\sqrt{\frac{2G}{qB_0 \left[1+\beta(\varepsilon z)\right]}}P_c\cos{g}\right)
\end{equation*}

\fi

\begin{gather*}
    \varepsilon R(\varepsilon;Q_c,P_c, g, G, z, P_z)=\frac{qB_0}{m}G\beta(\varepsilon z)-\frac{\varepsilon}{2m}P_zQ_cP_c \phi_1(\varepsilon z)-\varepsilon\sqrt{\frac{qB_0}{2}}\frac{P_zQ_c\sqrt{G}}{m}\sin{(g)} \phi_2(\varepsilon z)\,+\\
    +\,\varepsilon\frac{P_zP_c\sqrt{G}}{m\sqrt{2qB_0}}\cos{(g)} \phi_3(\varepsilon z)+\varepsilon^2\frac{Q_c^2P_c^2}{8m}\phi_4(\varepsilon z)+\varepsilon^2\frac{qB_0}{4m}Q_c^2G\sin^2{(g)}\phi_5(\varepsilon z)+\varepsilon^2\frac{P_cG}{4mqB_0}\cos^2{(g)}\phi_6(\varepsilon z)\,+\\
    +\,\varepsilon^2 \sqrt{\frac{qB_0}{2}}\frac{Q_c^2P_c\sqrt{G}}{2m}\sin{(g)}\phi_7(\varepsilon z)+\varepsilon^2\frac{Q_cP_c^2\sqrt{G}}{2m\sqrt{2qB_0}}\cos{(g)}\phi_8(\varepsilon z)+\varepsilon^2 \frac{Q_cP_c G}{2m}\sin{(g)}\cos{(g)}\phi_4(\varepsilon z)
\end{gather*}

\begin{align*}
    &\phi_1(\varepsilon z)=\frac{\beta'(\varepsilon z)}{1+\beta(\varepsilon z)} & &\phi_4(\varepsilon z)=\left[\frac{\beta'(\varepsilon z)}{1+\beta(\varepsilon z)}\right]^2 & &\phi_7(\varepsilon z)=\frac{\left[\beta'(\varepsilon z)\right]^2}{\left[1+\beta(\varepsilon z)\right]^{3/2}}\\
    &\phi_2(\varepsilon z)=\frac{\beta'(\varepsilon z)}{\sqrt{1+\beta(\varepsilon z)}}& &\phi_5(\varepsilon z)=\frac{\left[\beta'(\varepsilon z)\right]^2}{1+\beta(\varepsilon z)}& &\phi_8(\varepsilon z)=\frac{\left[\beta'(\varepsilon z)\right]^2}{\left[1+\beta(\varepsilon z)\right]^{5/2}}\\
    &\phi_3(\varepsilon z)=\frac{\beta'(\varepsilon z)}{\left[1+\beta(\varepsilon z)\right]^{3/2}} & &\phi_6(\varepsilon z)=\frac{\left[\beta'(\varepsilon z)\right]^2}{\left[1+\beta(\varepsilon z)\right]^3}
\end{align*}

We will now demonstrate how to evaluate the auxiliary Hamiltonian as an expansion in the perturbation parameter $\varepsilon$. Since the function $\beta$ can be defined to not have any constant term, we can write $\beta(\varepsilon z)=\sum_{j\geq1}\beta_j \varepsilon^j/j!$, where $\beta_j=\beta^{(j)}(0)$ (if we are Taylor-expanding about an arbitrary $z$-coordinate, we can always expand in terms of the displacement from the aforementioned). The unperturbed Hamiltonian is simply given by $\mathcal{K}(G,z,P_z)$, which we will call $H_0^0$, while the first-order perturbation is given by:

\begin{equation}
    \label{H-50}
    H_1^0(Q_c,P_c,g,G,z,P_z)=\frac{qB_0}{m}Gz\beta_1-\frac{P_z\beta_1}{2m}\left(Q_cP_c+\sqrt{2qB_0G}Q_c\sin{g}-\sqrt{\frac{2G}{qB_0}}P_c\cos{g}\right)
\end{equation} 

We now determine the auxiliary Hamiltonian $W_1$ and the new \textit{averaged} Hamiltonian by solving:

\begin{equation*}
    \left\{H_0^0,W_1\right\}=\overline{H}_1^0-H_1^0
\end{equation*}

where the Poisson bracket here is with respect to $(Q_c,P_c,g,G,z,P_z)$. However, the unperturbed Hamiltonian $H_0^0=H_0^0(G,z,P_z)$, so that all partial derivatives with respect to $g$, $Q_c$ and $P_c$ vanish: 

\begin{equation*}
    \left\{H_0^0,W_1\right\}=\left(\frac{\partial H_0^0}{\partial g}\frac{\partial W_1}{\partial G}-\frac{\partial H_0^0}{\partial G}\frac{\partial W_1}{\partial g}\right)+\left(\frac{\partial H_0^0}{\partial z}\frac{\partial W_1}{\partial P_z}-\frac{\partial H_0^0}{\partial P_z}\frac{\partial W_1}{\partial z}\right)+\left\{H_0^0, W_1\right\}_{Q_c, P_c}
\end{equation*}

\begin{equation}
    \label{H-51}
    \left\{H_0^0,W_1\right\}=-\frac{qB_0}{m}\frac{\partial W_1}{\partial g}-\frac{P_z}{m}\frac{\partial W_1}{\partial z}
\end{equation}

We now decompose the first-order correction of the exact Hamiltonian $H_1^0$ into its mean part $\overline{H}_1^0$, which should be independent of the gyroangle $g$, and its oscillatory part $\{H_1^0\}$. By comparison with (\ref{H-50}):

\begin{equation*}
    \overline{H}_1^0(Q_c,P_c,G,z,P_z)=\frac{qB_0}{m}Gz\beta_1-\frac{P_z\beta_1}{2m}Q_cP_c
\end{equation*}

Thus, we search for $W_1(Q_c,P_c,g,g,z,P_z)$ such that:

\begin{equation*}
    \frac{qB_0}{m}\frac{\partial W_1}{\partial g}+\frac{P_z}{m}\frac{\partial W_1}{\partial z}=\frac{\beta_1}{m}\sqrt{\frac{qB_0}{2}}P_zQ_c\sqrt{G}\sin{g}-\frac{\beta_1}{m\sqrt{2qB_0}}P_zP_c\sqrt{G}\cos{g}
\end{equation*}

Since the right-hand side of the above equation has no $z$-dependence, including a $z$-coordinate contribution would pertain to a mere constant, so that we may as well set $\partial W_1/\partial z=0$. Simply integrating:

\begin{equation}
    W_1(Q_c,P_c,g,G,P_z)=-\frac{\beta_1}{\sqrt{2qB_0}}P_zQ_c\sqrt{G}\cos{g}-\frac{\beta_1}{\sqrt{2q^3B_0^3}}P_zP_c\sqrt{G}\sin{g}
\end{equation}








































\newpage

\section{Constant uniform electric field}

The equation of motion is, in this case, trivial to solve:

\begin{equation}
    \diff{\mathbf{p}}{t}=q\mathbf{E}_0\implies\mathbf{p}(t)=\mathbf{p}_0+q\mathbf{E}_0t
\end{equation}

\begin{equation}
    \mathbf{v}(t)=\frac{\mathbf{p}(t)}{\gamma(t)m}=\frac{\mathbf{p}_0+q\mathbf{E}_0t}{\sqrt{1+\frac{(\mathbf{p}_0+q\mathbf{E}_0t)^2}{m^2c^2}}}\quad\quad\lim_{t\to\infty}\|\mathbf{v}(t)\|=c
\end{equation}

\section{Constant uniform magnetic field}

\label{sec2-2}

The equation of motion is also trivial to solve, as in the previous case:

\begin{equation}
    \diff{\mathbf{p}}{t}=q\mathbf{v}\times\mathbf{B}_0\implies\diff{\mathbf{v}}{t}=\mathbf{v}\times\frac{q\mathbf{B_0}}{\gamma m}\equiv\mathbf{v}\times\mathbf{\Omega_c}\quad\quad\mathbf{\Omega}_c=\frac{q\mathbf{B}_0}{\gamma m}
\end{equation}

where $\Omega_c$ designates the relativistic angular velocity of the particle considered. The norm of the charge's velocity $\|\mathbf{v}\|$ is conserved:

\begin{equation}
    \mathbf{v}\cdot\diff{\mathbf{v}}{t}=\frac{1}{2}\diff{}{t}\left(\mathbf{v}^2\right)=\diff{}{t}\left(\|\mathbf{v}\|\right)=\mathbf{v}\cdot\left(\mathbf{v}\times\Omega_c\right)=0\implies\diff{\gamma}{t}=0
\end{equation}

The motion along the magnetic field is uniform. Let us choose $\mathbf{B}_0$ along the $\mathbf{\hat{z}}$ axis, then:

\begin{equation}
    \diff{}{t}\left(\mathbf{v}\cdot\mathbf{\hat{z}}\right)=0\quad\quad\mathbf{v}-(\mathbf{v}\cdot\mathbf{\hat{z}})\,\mathbf{\hat{z}}=v_\mathrm{x}\,\mathbf{\hat{x}}+v_\mathrm{y}\,\mathbf{\hat{y}}\implies v_\perp=\sqrt{v_\mathrm{x}^2+v_\mathrm{y}^2}
\end{equation}

\begin{equation}
    \diff{v_\mathrm{x}}{t}=\frac{qB_0}{\gamma m}v_\mathrm{y}\qquad\diff{v_\mathrm{y}}{t}=-\frac{qB_0}{\gamma m}v_\mathrm{x}
\end{equation}

\begin{equation}
    \frac{\mathrm{d}^2v_\mathrm{x}}{\mathrm{d}t^2}=\frac{qB_0}{\gamma m}\diff{v_\mathrm{y}}{t}=-\left(\frac{qB_0}{\gamma m}\right)^2v_\mathrm{x}\qquad\frac{\mathrm{d}^2v_\mathrm{y}}{\mathrm{d}t^2}=-\frac{qB_0}{\gamma m}\diff{v_\mathrm{x}}{t}=-\left(\frac{qB_0}{\gamma m}\right)^2v_\mathrm{y}
\end{equation}

\begin{equation}
    v_\mathrm{x}(t)=v_\perp\cos(\Omega_ct+\varphi)\qquad v_\mathrm{y}(t)=v_\perp\sin(\Omega_ct+\varphi)
\end{equation}

\begin{equation}
    x(t)=\int v_\mathrm{x}(t)\,\mathrm{d}t=x_0+\frac{v_\perp}{\Omega_c}\sin(\varphi)+\frac{v_\perp}{\Omega_c}\sin(\Omega_ct+\varphi)
\end{equation}

\begin{equation}
    y(t)=\int v_\mathrm{y}(t)\,\mathrm{d}t=y_0-\frac{v_\perp}{\Omega_c}\cos(\varphi)-\frac{v_\perp}{\Omega_c}\cos(\Omega_ct+\varphi)
\end{equation}

We will take further advantage of this by defining the angle $\alpha=\varphi-\frac{\pi}{2}$:

\begin{equation}
    x(t)=x_\text{g}+r_\text{L}\cos(\Omega_ct+\alpha)\quad\quad\quad y(t)=y_\text{g}+r_\text{L}\sin(\Omega_ct+\alpha)
\end{equation}

\begin{equation}
    x_\text{g}=x_0-r_\text{L}\cos(\alpha)\qquad y_\text{g}=y_0-r_\text{L}\sin(\alpha)\qquad r_\text{L}=\frac{\gamma mv_\perp}{\|q\mathbf{B_0}\|}=\frac{\|\mathbf{p}_\perp\|}{\|q\mathbf{B}_0\|}
\end{equation}

The centre of the particle rotation orbits is termed the \textit{guiding centre} or \textit{gyro-centre}. It is defined by the vector $\mathbf{r}_\text{g}(t)=x_\text{g}\,\mathbf{\hat{x}}+y_\text{g}\,\mathbf{\hat{y}}+z(t)\,\mathbf{\hat{z}}$. The quantity $r_\text{L}$ is known as the \textit{Larmor radius} or \textit{gyro radius}. It should be noted that the direction of rotation for positively charged particles is clockwise (left-handed), whereas the direction of rotation for negatively charged particles is counterclockwise (right-handed). This results in the magnetic dipole created by the motion of the concerned particle being oriented in the opposite direction to the external magnetic field. In other words, the particles tend to reduce the magnetic field imposed on them. This is precisely the property of magnetic materials termed \textit{diamagnetism}. Plasma is therefore a diamagnetic medium.

\section{Constant uniform magnetic and electric fields}

The non-relativistic Lorentz equation appears in the following form:

\begin{equation}
    m\diff{\mathbf{v}}{t}=q\left(\mathbf{E}+\mathbf{v}\times\mathbf{B}\right)
\end{equation}

We will decompose the velocity $\mathbf{v}$ into three components: the velocity $\mathbf{v}_\parallel$ parallel to the magnetic field $\mathbf{B}$, the gyration velocity $\mathbf{v}_\text{L}$, and another velocity, introduced under the term \textit{drift velocity}, $\mathbf{v}_\mathrm{d}$:

\begin{equation}
    m\,\diff{}{t}\left(\mathbf{v}_\parallel+\mathbf{v}_\mathrm{d}+\mathbf{v}_\text{L}\right)=q\left[\mathbf{E}+\left(\mathbf{v}_\parallel+\mathbf{v}_\mathrm{d}+\mathbf{v}_\text{L}\right)\times\mathbf{B}\right]
\end{equation}

It is easy to see that one must define the drift velocity such that:

\begin{equation}
    \label{eq2-17}
    \mathbf{E}+\mathbf{v}_\mathrm{d}\times\mathbf{B}=\mathbf{0}
\end{equation}

so that we have a simple cyclotron motion satisfied by $\mathbf{v}_\text{L}$ as in section (\ref{sec2-2}). By taking the cross product of equation (\ref{eq2-17}) with $\mathbf{B}$, it follows:

\begin{equation}
    \mathbf{E}\times\mathbf{B}+\left(\mathbf{v}_\mathrm{d}\times\mathbf{B}\right)\times\mathbf{B}=\mathbf{E}\times\mathbf{B}+(\underbrace{\mathbf{v}_\mathrm{d}\cdot\mathbf{B}}_\mathbf{0})\mathbf{B}-\mathbf{B}^2\mathbf{v}_\mathrm{d}=\mathbf{0}
\end{equation}

The expression for the drift velocity is thus easily determined:

\begin{equation}
    \mathbf{v}_\mathrm{d}=\frac{\mathbf{E}\times\mathbf{B}}{\mathbf{B}^2}
\end{equation}

It should be noted that the drift velocity $\mathbf{v}_\mathrm{d}$ is independent of the particle's properties; thus, it points in the same direction for positive and negative charges. The physical intuition behind the drift velocity is that, in the frame moving at this velocity, the electric field is transformed so that it is zero. When introducing relativistic effects, one must distinguish between the following two cases: $\|\mathbf{E}\|<c\|\mathbf{B}\|$ and $\|\mathbf{E}\|>c\|\mathbf{B}\|$. Otherwise, we would encounter contradictions. Without loss of generality, let $\mathbf{E}\perp\mathbf{B}$, then:

\begin{equation}
    \|\mathbf{E}\|>c\|\mathbf{B}\|\implies\|\mathbf{v}_\mathrm{d}\|=\Bigg\|\frac{\mathbf{E}\times\mathbf{B}}{\mathbf{B}^2}\Bigg\|>c\quad\perp
\end{equation}

Let us consider the frame moving, relative to the laboratory frame where electric and magnetic fields are defined such that $\|\mathbf{E}\|<c\|\mathbf{B}\|$, at the following velocity:

\begin{equation}
    \mathbf{v}_\mathrm{d}=\frac{\mathbf{E}\times\mathbf{B}}{\|\mathbf{B}\|^2}\implies\|\mathbf{v}_\mathrm{d}\|<c
\end{equation}

Using Lorentz transformations, we recover the following fields in the new frame:

\begin{equation}
    \mathbf{E}'=\gamma_\mathrm{d}\left(\mathbf{E}+\mathbf{v}_\mathrm{d}\times\mathbf{B}\right)-(\gamma_\mathrm{d}-1)\big(\underbrace{\mathbf{E}\cdot\mathbf{\hat{v}}_\mathrm{d}}_{0}\big)\,\mathbf{\hat{v}}_\mathrm{d}=\gamma_\mathrm{d}\left[\mathbf{E}+\frac{(\mathbf{E}\times\mathbf{B})}{\|\mathbf{B}\|^2}\times\mathbf{B}\right]=\gamma_\mathrm{d}\frac{\mathbf{E}\cdot\mathbf{B}}{\|\mathbf{B}\|^2}\mathbf{B}
\end{equation}

\begin{equation}
    \mathbf{B}'=\gamma_\mathrm{d}\left(\mathbf{B}-\frac{\mathbf{v}_\mathrm{d}\times\mathbf{E}}{c^2}\right)-(\gamma_\mathrm{d}-1)(\underbrace{\mathbf{B}\cdot\mathbf{\hat{v}_d}}_0)\,\mathbf{\hat{v}}_d=\gamma_\mathrm{d}\left[\mathbf{B}+\frac{(\mathbf{E}\cdot\mathbf{B})\mathbf{E}-\|\mathbf{E}\|^2\mathbf{B}}{c^2\|\mathbf{B}\|^2}\right]
\end{equation}

In the specific case where $\mathbf{E}\perp\mathbf{B}$, we observe a significant simplification of the preceding expressions:

\begin{equation}
    \mathbf{E}'=0\quad\text{et}\quad\mathbf{B}'=\gamma_\mathrm{d}\left(1-\frac{\mathbf{E}^2}{c^2\mathbf{B}^2}\right)\mathbf{B}=\frac{\mathbf{B}}{\gamma_\mathrm{d}}
\end{equation}

Similarly, in the case where $\|\mathbf{E}\|>c\|\mathbf{B}\|$, and furthermore $\mathbf{E}\perp\mathbf{B}$, we simply admit that the expression for the drift velocity is of the following form:

\begin{equation}
    \mathbf{v}_\mathrm{d}=c^2\frac{\mathbf{E}\times\mathbf{B}}{\mathbf{E}^2}\implies\mathbf{E}'=\frac{\mathbf{E}}{\gamma_\mathrm{d}}\quad\text{et}\quad\mathbf{B}'=\mathbf{0}
\end{equation}

\section{Constant uniform magnetic field and force field}

The equation of motion in the case where the considered particle undergoes a force $\mathbf{F}$ reads:

\begin{equation}
    m\diff{\mathbf{v}}{t}=\mathbf{F}+q\mathbf{v}\times\mathbf{B}=q\left(\frac{\mathbf{F}}{q}+\mathbf{v}\times\mathbf{B}\right)
\end{equation}

Identically to the previous section, we can prove that the drift velocity appears in the following form, the derivation of which is omitted due to redundancy:

\begin{equation}
    \label{eq2-27}
    \mathbf{v}_\mathrm{d}=\frac{1}{q}\frac{\mathbf{F}\times\mathbf{B}}{\mathbf{B}^2}
\end{equation}

\section{Constant non-uniform magnetic field}

\label{sec2-5}

\subsection{Direction of variation perpendicular to the field}

\label{sec2-5-1}

We examine the motion of a particle of charge $q$ and mass $m$ subjected to a magnetic field $\mathbf{B}$ characterized by a non-zero dyadic tensor $\nabla\mathbf{B}$. For the sake of simplicity and physical intuition, we will adopt the following assumptions:

\begin{enumerate}
    \item\label{hyp2-5-1-1} Any velocity considered is much smaller than the speed of light. This allows us to neglect relativistic effects, which are indeed, for the most part, negligible.
    \item\label{hyp2-5-1-2} The magnetic field undergoes minuscule variations on the scale of the Larmor radius of the particle considered, i.e. $r_\text{L}\|\nabla\mathbf{B}\|\ll\|\mathbf{B}\|$, the approximation $\mathbf{B}\simeq\mathbf{B}_0+(\mathbf{r}\cdot\nabla)\mathbf{B}$ holds very well.
    \item\label{hyp2-5-1-3} The dyadic tensor is a tensor with a single non-zero component, $\mathbf{\hat{x}}\mathbf{\hat{B}}$, where $\mathbf{\hat{x}}\perp\mathbf{B}$.
    \item\label{hyp2-5-1-4} The drift velocity is much smaller than the Larmor velocity of the particle considered, i.e. $\|\mathbf{v}_\mathrm{d}\|\ll\|\mathbf{v}_\text{L}\|$ holds very well. The validity of this assumption will be examined shortly.
    \item\label{hyp2-5-1-5} The drift velocity indeed undergoes variations, but these variations are perfectly negligible. The validity of this assumption will be established at the close of the subsection.
\end{enumerate}

By virtue of the preceding analyses, we impose the velocity decomposition $\mathbf{v}=\mathbf{v}_\mathrm{d}+\mathbf{v}_\text{L}$ such that the Lorentz equation applied to the particle considered reads:

\begin{equation}
    m\diff{\mathbf{v}}{t}\equiv m\diff{\mathbf{v}_\text{L}}{t}=q\left(\mathbf{v}_\mathrm{d}+\mathbf{v}_\text{L}\right)\times\mathbf{B}
\end{equation}

Using assumption (\ref{hyp2-5-1-2}), we expand the cross product as follows:

\begin{equation}
    m\diff{\mathbf{v}_\text{L}}{t}=q\bigg[\mathbf{v}_\text{L}\times\mathbf{B}_0+\underbrace{\mathbf{v}_\mathrm{d}\times\mathbf{B}_0+(\mathbf{v}_\mathrm{d}+\mathbf{v}_\text{L})\times(\mathbf{r}\cdot\nabla)\mathbf{B}}_\mathbf{0}\bigg]
\end{equation}

We impose the drift-velocity condition, such that the underlined sum results in zero:

\begin{equation}
    \label{eq2-30}
    \mathbf{v}_\mathrm{d}\times\mathbf{B}_0+\mathbf{v}_\mathrm{d}\times(\mathbf{r}\cdot\nabla)\mathbf{B}+\underbrace{\mathbf{v}_\text{L}\times(\mathbf{r}\cdot\nabla)\mathbf{B}}_{\langle\cdots\rangle=\mathbf{0}}=\mathbf{0}
\end{equation}

We identify the second term as the most negligible thanks to assumptions (\ref{hyp2-5-1-2}) and (\ref{hyp2-5-1-4}). Before beginning the resolution of the condition given by equation (\ref{eq2-30}), we impose the decomposition of the position vector $\mathbf{r}=\mathbf{r}_\text{g}+\mathbf{r}_\text{L}$, in other words, into guiding centre and Larmor radius:

\begin{equation}
    \label{eq2-31}
    \mathbf{v}_\mathrm{d}\times\mathbf{B}+\mathbf{v}_\text{L}\times(\mathbf{r}_\text{L}\cdot\nabla)\mathbf{B}+\mathbf{v}_\text{L}\times(\mathbf{r}_\text{g}\cdot\nabla)\mathbf{B}=\mathbf{0}
\end{equation}

We will proceed by taking the average of equation (\ref{eq2-31}) over a rotation period. It is easy to see that the last term of the previous equation vanishes when its average is taken. Assumption (\ref{hyp2-5-1-3}) gives:

\begin{equation}
    \mathbf{r}_\text{L}=\frac{v_{\perp\mathbf{B}}}{\Omega_\text{c}}\begin{pmatrix}
        \cos(\Omega_\text{c}t+\varphi_0)\\
        \sin(\Omega_\text{c}t+\varphi_0)
    \end{pmatrix}\implies\mathbf{v}_\text{L}=v_{\perp\mathbf{B}}\begin{pmatrix}
        -\sin(\Omega_\text{c}t+\varphi_0)\\
        \cos(\Omega_\text{c}t+\varphi_0)
    \end{pmatrix}\quad\text{et}\quad(\mathbf{r}\cdot\nabla)\mathbf{B}=x\frac{\mathrm{d}\mathbf{B}}{\mathrm{d}x}
\end{equation}

We examine the second term of equation (\ref{eq2-31}) more closely by decomposing it along $\mathbf{\hat{x}}$ and $\mathbf{\hat{y}}\perp\mathbf{B}$:

\begin{equation}
    \left[\mathbf{v}_\text{L}\times\left(x\frac{\mathrm{d}\mathbf{B}}{\mathrm{d}x}\right)\right]\cdot\mathbf{\hat{x}}=v_\text{Ly}x\frac{\mathrm{d}B}{\mathrm{d}x}\quad\text{et}\quad\left[\mathbf{v}_\text{L}\times\left(x\frac{\mathrm{d}\mathbf{B}}{\mathrm{d}x}\right)\right]\cdot\mathbf{\hat{y}}=-v_\text{Lx}x\frac{\mathrm{d}B}{\mathrm{d}x}
\end{equation}

\begin{equation}
    \langle v_\text{Ly}x\rangle=\frac{v_{\perp\mathbf{B}}^2}{\Omega_c}\left\langle\cos(\Omega_ct+\varphi_0)\cos(\Omega_ct+\varphi_0)\right\rangle=\frac{v_{\perp\mathbf{B}}^2}{2\Omega_c}
\end{equation}

\begin{equation}
    \langle v_\text{Lx}x\rangle=\frac{v_{\perp\mathbf{B}}}{\Omega_c}\left\langle-\sin(\Omega_ct+\varphi_0)\cos(\Omega_ct+\varphi_0)\right\rangle=0
\end{equation}

Returning to equation (\ref{eq2-31}), we obtain the expression for the drift velocity in an analogous manner as in the previous sections, hence:

\begin{equation}
    \label{eq2-36}
    \mathbf{v}_{\nabla\mathbf{B}}=\frac{1}{q}\frac{mv_{\perp}^2}{2}\frac{\mathbf{B}\times(\mathbf{B}\cdot\nabla)\mathbf{B}}{\mathbf{B}^4}
\end{equation}

This drift velocity is commonly referred to as \textit{grad B-drift}. It is worth confirming the consistency between the assumptions established at the beginning of the discussion and the result obtained. We first test assumption (\ref{hyp2-5-1-4}) and find that it is well verified. Let us note assumption (\ref{hyp2-5-1-3}):

\begin{equation}
    \|\mathbf{v}_{\nabla\mathbf{B}}\|=\Bigg\|\frac{1}{q}\frac{mv_\perp^2}{2}\frac{\mathbf{B}\times(\mathbf{B}\cdot\nabla)\mathbf{B}}{\mathbf{B}^4}\Bigg\|=\frac{mv_\perp^2}{2|q|r_\text{L}}\frac{r_\text{L}\|\nabla\mathbf{B}\|}{\mathbf{B}^2}=\frac{mv_\perp^2}{2|q|}\frac{\|q\mathbf{B}\|}{mv_\perp}\frac{r_\text{L}\|\nabla\mathbf{B}\|}{\mathbf{B}^2}\ll\frac{v_\perp}{2}=\frac{\|\mathbf{v}_\text{L}\|}{2}
\end{equation}

The variation of the drift velocity is given by the following:

\begin{equation}
    \delta\|\mathbf{v}_{\nabla\mathbf{B}}\|=\frac{mv_\perp^2}{2|q|}\delta\left(\frac{\|\nabla\mathbf{B}\|}{\mathbf{B}^2}\right)
\end{equation}

The order of the variation of the perpendicular velocity is the same as that of the magnetic field. Thus, assumption (\ref{hyp2-5-1-2}) ensures the minimality of the variation of $\mathbf{v}_{\nabla\mathbf{B}}$.

\subsection{Direction of variation parallel to the field}

\label{sec2-5-2}

When a magnetic field experiences intensification due to the dyadic tensor $\nabla\mathbf{B}$, whose component $\mathbf{\hat{z}}\mathbf{\hat{B}}$ is non-zero, this has the effect of compiling the field lines alongside each other. In other words, the magnetic field lines locally form a section cut by a cone, and the radial component is no longer zero, i.e. $\mathbf{B}\cdot\mathbf{\hat{r}}\neq0$, however, such that the inequality $\mathbf{B}\cdot\mathbf{\hat{r}}\ll\|\mathbf{B}\|$ is satisfied. This strong inequality is analogous to assumption (\ref{hyp2-5-1-2}) of subsection (\ref{sec2-5-1}). This entails that the magnetic field varies very little along the $\mathbf{\hat{z}}$ axis.
By taking the average of the Lorentz force over a rotation period, there exists a non-zero component of this force parallel to the direction of magnetic field increase $\mathbf{\hat{z}}$:

\begin{equation}
    \langle\mathbf{F}\cdot\mathbf{\hat{z}}\rangle=\langle q(\mathbf{v}\times\mathbf{B})\cdot\mathbf{z}\rangle=q\langle\mathbf{v}\cdot(\mathbf{B}\times\mathbf{\hat{z}})\rangle=qv_\perp\mathbf{B}\cdot\mathbf{\hat{r}}=-qv_\perp\|\mathbf{B}\|\sin(\theta)
\end{equation}

where $\theta$ is the angle between the magnetic field $\mathbf{B}$ and the $\mathbf{\hat{z}}$ axis in the trigonometric sense. Using a Maxwell equation, we can relate the components of the magnetic field along the axes $\mathbf{\hat{r}}$ and $\mathbf{\hat{z}}$:

\begin{equation}
    \nabla\cdot\mathbf{B}=0\stackrel{B_\varphi=0}{\implies}\frac{1}{r}\frac{\partial}{\partial r}(rB_\text{r})+\frac{\partial B_\text{z}}{\partial z}=0\implies rB_\text{r}=-\int r\frac{\partial B_\text{z}}{\partial z}\,\mathrm{d}r
\end{equation}

Under the assumption that the Larmor radius is sufficiently small for the variation $\partial_zB_\text{z}$ to be approximately constant, we obtain:

\begin{equation}
    rB_r\,\bigg|_0^{r_\text{L}}\simeq-\frac{\partial B_\text{z}}{\partial z}\int_0^{r_\text{L}}r\,\mathrm{d}r=-\frac{1}{2}r_\text{L}^2\frac{\partial B_\text{z}}{\partial z}\implies B_\text{r}\simeq-\frac{r_\text{L}}{2}\frac{\partial B_\text{z}}{\partial z}\implies\sin(\theta)=\frac{r_\text{L}}{2}\frac{1}{\|\mathbf{B}\|}\frac{\partial B_\text{z}}{\partial z}
\end{equation}

Consequently, the particle undergoes a recoil force given by the following expression:

\begin{equation}
    \label{eq2-42}
    \langle\mathbf{F}\cdot\mathbf{\hat{z}}\rangle=-\frac{qv_\perp r_\text{L}}{2}\frac{\partial B_\text{z}}{\partial z}=-\frac{mv_\perp^2}{2\|\mathbf{B}\|}\frac{\partial B_\text{z}}{\partial z}
\end{equation}

\section{Magnetic field presenting a radius of curvature}

\label{sec2-6}

We will next analyse the motion of a particle of charge $q$ and mass $m$ in a uniform and constant magnetic field of magnitude $\|\mathbf{B}\|$ and radius of curvature $\mathbf{R}$. By entering the rotating frame in which the particle follows a simple cyclotron motion, it undergoes inertial forces, in particular the centrifugal force and the Coriolis force. For simplicity, we admit straight away that the Coriolis effect does not influence the particle's motion significantly at all. The proof of this assertion relies on the fact that the Coriolis effect vanishes when taking its average over a rotation period. Let $\mathbf{v}_\parallel=(\mathbf{v}\cdot\mathbf{\hat{B}})\,\mathbf{\hat{B}}$ be the component of the particle's velocity parallel to the magnetic field, assumed constant.
Then, the centrifugal force reads:

\begin{equation}
    \mathbf{F}_\text{cf}=mv_\parallel^2\frac{\mathbf{R}}{\mathbf{R}^2}
\end{equation}

Using equation (\ref{eq2-27}), the drift velocity induced by the curvature of the magnetic field is:

\begin{equation}
    \mathbf{v}_\mathbf{R}=\frac{mv_\parallel^2}{q}\frac{\mathbf{R}\times\mathbf{B}}{\mathbf{R}^2\mathbf{B}^2}
\end{equation}

The fact that the magnitude of the magnetic field is constant does not prevent its gradient from being non-zero. It is therefore appropriate to relate the dyadic tensor $\nabla\mathbf{B}$ to the radius of curvature. By introducing a cylindrical coordinate system such that $\mathbf{B}\cdot\mathbf{\hat{r}}=0=\mathbf{B}\cdot\mathbf{\hat{z}}$, and by taking advantage of a Maxwell equation, we will achieve this elegantly:

\begin{equation}
    \nabla\times\mathbf{B}=\mathbf{0}\implies(\nabla\times\mathbf{B})\cdot\mathbf{\hat{z}}=\frac{1}{r}\frac{\partial}{\partial r}(rB_\varphi)-\frac{1}{r}\underbrace{\frac{\partial B_r}{\partial\varphi}}_0=0
\end{equation}

\begin{equation}
    \frac{\partial B_\varphi}{\partial r}=-\frac{B_\varphi}{r}\implies\nabla\mathbf{B}=-\frac{\mathbf{R}\mathbf{B}}{\mathbf{R}^2}\implies\left(\mathbf{B}\cdot\nabla\right)\mathbf{B}\perp\mathbf{B}
\end{equation}

Using the result (\ref{eq2-36}), the drift velocity induced by the magnetic field gradient, in this specific case, reads:

\begin{equation}
    \mathbf{v}_{\nabla\mathbf{B}}=\frac{1}{q}\frac{mv_\perp^2}{2}\frac{\mathbf{B}}{\|\mathbf{B}\|^3}\times\left(-\frac{\mathbf{R}}{\mathbf{R}^2}\right)\|\mathbf{B}\|=\frac{1}{q}\frac{mv_\perp^2}{2}\frac{\mathbf{R}\times\mathbf{B}}{\mathbf{R}^2\mathbf{B}^2}
\end{equation}

The total drift is the simple sum of the two drift velocities obtained:

\begin{equation}
    \label{eq2-48}
    \mathbf{v}_\mathrm{d}=\mathbf{v}_{\nabla\mathbf{B}}+\mathbf{v}_\mathbf{R}=\frac{1}{q}\left(mv_\parallel^2+\frac{1}{2}mv_\perp^2\right)\frac{\mathbf{R}\times\mathbf{B}}{\mathbf{R}^2\mathbf{B}^2}
\end{equation}

It should be noted that the average of the preceding drift velocity can be evaluated quickly thanks to the
equipartition theorem:

\begin{equation}
    \left\langle\frac{1}{2}mv_\parallel^2\right\rangle=\frac{1}{2}k_\text{B}T\quad\text{et}\quad\left\langle\frac{1}{2}mv_\perp^2\right\rangle=2\times\frac{1}{2}k_\text{B}T\implies\langle\mathbf{v}_\mathrm{d}\rangle=\frac{2k_\text{B}T}{q}\frac{\mathbf{R}\times\mathbf{B}}{\mathbf{R}^2\mathbf{B}^2}
\end{equation}

\section{Variable electric and constant magnetic field}

The drift velocity undergoes a temporal variation because the electric field, for its part, depends on time:

\begin{equation}
    \mathbf{v}_\mathrm{d}(t)=\frac{\mathbf{E}(t)\times\mathbf{B}}{\mathbf{B}^2}\implies\mathbf{a}_\mathrm{d}=\diff{\mathbf{v}_\mathrm{d}}{t}=\diff{}{t}\left(\frac{\mathbf{E}\times\mathbf{B}}{\mathbf{B}^2}\right)
\end{equation}

The guiding centre undergoes acceleration, as is also the case for the particle considered. We will examine the specific case of an oscillatory electric field perpendicular to the magnetic field:

\begin{equation}
    m\diff{\mathbf{v}}{t}=q\left(\mathbf{E}e^{i\omega t}+\mathbf{v}\times\mathbf{B}\right)
\end{equation}

Complex notation implies only the real component of the number $e^{\mathrm{i}\omega t}$. Multiplication by $i\omega$ is thus equivalent to the derivative $\mathrm{d}/\mathrm{d}t$. We impose the following decomposition:

\begin{equation}
    \label{eq2-52}
    \mathbf{v}=\mathbf{v}_\mathrm{d}e^{i\omega t}+\mathbf{v}_\text{L}\quad\text{où}\quad m\diff{\mathbf{v}_\text{L}}{t}=q\mathbf{v}_\text{L}\times\mathbf{B}
\end{equation}

\begin{equation}
    \label{eq2-53}
    im\omega\,\mathbf{v}_\mathrm{d}=q\left(\mathbf{E}+\mathbf{v}_\mathrm{d}\times\mathbf{B}\right)e^{i\omega t}\implies im\omega\,\mathbf{v}_\mathrm{d}\times\mathbf{B}=q\left[\mathbf{E}\times\mathbf{B}+(\mathbf{B\cdot\mathbf{v}_\mathrm{d}})\mathbf{B}-\mathbf{B}^2\mathbf{v}_\mathrm{d}\right]e^{i\omega t}
\end{equation}

To eliminate $\mathbf{v}_\mathrm{d}\times\mathbf{B}$, we multiply the first equation of (\ref{eq2-53}) by $-im\omega$ and the second by $q$:

\begin{equation}
    m^2\omega^2\mathbf{v}_\mathrm{d}=-\mathrm{i}mq\omega\left(\mathbf{E}+\mathbf{v}_\mathrm{d}\times\mathbf{B}\right)e^{\mathrm{i}\omega t}\quad\text{et}\quad\mathrm{i}mq\omega\,\mathbf{v}_\mathrm{d}\times\mathbf{B}=q^2\left[\mathbf{E}\times\mathbf{B}+(\mathbf{B}\cdot\mathbf{v}_\mathrm{d})\mathbf{B}-\mathbf{B}^2\mathbf{v}_\mathrm{d}\right]e^{\mathrm{i}\omega t}
\end{equation}

\begin{equation}
    m^2\omega^2\mathbf{v}_\mathrm{d}+q^2\left(\mathbf{E}\times\mathbf{B}-\mathbf{B}^2\mathbf{v}_\mathrm{d}\right)=\mathrm{i}mq\omega\,\mathbf{E}\equiv mq\diff{\mathbf{E}}{t}
\end{equation}

This drift velocity is also known as \textit{polarization drift}:

\begin{equation}
    \mathbf{v}_\mathrm{d}=\left(\frac{\mathbf{E}\times\mathbf{B}}{\mathbf{B}^2}+\frac{m}{q\mathbf{B}^2}\diff{\mathbf{E}}{t}\right)\Bigg/\left(1-\frac{m^2\omega^2}{q^2\mathbf{B}^2}\right)
\end{equation}

We are obliged to follow a different procedure when these temporal variations of the electric field do not present oscillations.

\section{Non-uniform electric and uniform magnetic field}

The electric field is assumed to be slightly non-uniform, analogously to the magnetic field, in section (\ref{sec2-5}). In other words, we adopt the following assumption:

\begin{equation}
    \mathbf{E}\simeq\mathbf{E}_0+(\mathbf{r}\cdot\nabla)\mathbf{E}\quad\text{where }\|\mathbf{r}\|\text{ is of the same order as the Larmor radius}
\end{equation}

\begin{equation}
    m\diff{\mathbf{v}}{t}=q\left[\mathbf{E}(r)+\mathbf{v}\times\mathbf{B}\right]
\end{equation}

We seek a solution for $\mathbf{v}_\mathrm{d}$ such that the following equation is satisfied:

\begin{equation}
    m\left\langle\diff{\mathbf{v}_\mathrm{d}}{t}\right\rangle=q\left\langle\mathbf{E}(r)+\mathbf{v}_\mathrm{d}\times\mathbf{B}\right\rangle=q\left[\langle\mathbf{E}\rangle+\mathbf{v}_\mathrm{d}\times\mathbf{B}\right]=\mathbf{0}
\end{equation}

We know the solution very well, so we only need to determine the average magnetic field over a rotation period. The averages of the mixed terms are all zero thanks to orthogonality, so we are left with:

\begin{equation}
    \langle\mathbf{E}\rangle=\mathbf{E}_0+(\langle\mathbf{r}_\text{L}\rangle\cdot\nabla)\mathbf{E}+\frac{1}{2}\langle(\mathbf{r}_\text{L}\cdot\mathbf{\hat{x}})^2\rangle\frac{\partial^2\mathbf{E}}{\partial x^2}+\frac{1}{2}\langle(\mathbf{r}_\text{L}\cdot\mathbf{\hat{y}})^2\rangle\frac{\partial^2\mathbf{E}}{\partial y^2}
\end{equation}

We easily see $\langle\mathbf{r}_\text{L}\rangle=\mathbf{0}$, $\langle(\mathbf{r}_\text{L}\cdot\mathbf{\hat{x}})^2\rangle=r_\text{L}^2/2$ and likewise for the component along $\mathbf{\hat{y}}$:

\begin{equation}
    \mathbf{v}_\mathrm{d}=\mathbf{v}_\mathbf{E}+\mathbf{v}_{\nabla\mathbf{E}}=\frac{\mathbf{E}\times\mathbf{B}}{\mathbf{B}^2}+\frac{r_\text{L}^2}{4}\nabla^2\mathbf{E}=\left(1+\frac{1}{4}\frac{m^2v_\perp^2}{q^2\mathbf{B}^2}\nabla^2\right)\frac{\mathbf{E}\times\mathbf{B}}{\mathbf{B}^2}
\end{equation}

\section{Adiabatic invariants}

Most periodic physical systems are characterized by a physical quantity which, despite possible \textit{slow} deformations of certain parameters, remains constant. This quantity is commonly designated by the term \textit{adiabatic invariant}. More precisely, adiabatic invariants are first-order approximations of \textit{Poincaré invariants}. We will adopt an analytical mechanics approach in order to study and theoretically develop the basic notions. Subsequently, we will draw on the preceding sections to illustrate several examples of adiabatic invariants. A Poincaré invariant takes the following form:

\begin{equation}
    \mathcal{I}=\oint_{\mathcal{C}(t)}\mathbf{p}\cdot\mathrm{d}\mathbf{q}
\end{equation}

where all points belonging to the closed curve $\mathcal{C}(t)$ in phase space move according to the system's equations of motion. We introduce a periodic variable $s\in\mathcal{S}$, which parametrizes the points of the curve $\mathcal{C}(t)$. The general coordinates of a point on this curve are then written in the form $q_i = q_i(s,t)$ and $p_i = p_i(s,t)$. The time derivative then reads:

\begin{equation}
    \diff{\mathcal{I}}{t}=\oint_{\mathcal{S}}\left(p_i\frac{\partial^2q_i}{\partial t\partial s}+\frac{\partial p_i}{\partial t}\frac{\partial q_i}{\partial s}\right)\mathrm{d}s
\end{equation}

We integrate the first term by parts and then substitute Hamilton's equations:

\begin{equation}
    \diff{\mathcal{I}}{t}=\oint_\mathcal{S}\left(-\frac{\partial q_i}{\partial t}\frac{\partial p_i}{\partial s}+\frac{\partial p_i}{\partial t}\frac{\partial q_i}{\partial s}\right)\mathrm{d}s=-\oint_\mathcal{S}\left(\frac{\partial\mathcal{H}}{\partial p_i}\frac{\partial p_i}{\partial s}+\frac{\partial\mathcal{H}}{\partial q_i}\frac{\partial q_i}{\partial s}\right)\mathrm{d}s
\end{equation}

where $\mathcal{H}(\mathbf{q},\mathbf{p},t)$ represents the Hamiltonian of the system. If the system is conservative, i.e., if the Hamiltonian doesn't explicitly depend on time, the integrand is the total derivative of the Hamiltonian with respect to the variable $s$ along the curve $\mathcal{C}$. The time-independent Hamiltonian being a function of state, the integral evaluates to zero, thus demonstrating that $\mathcal{I}$ is indeed invariant:

\begin{equation}
    \diff{\mathcal{I}}{t}=-\oint_\mathcal{S}\diff{\mathcal{H}}{s}\,\mathrm{d}s=0
\end{equation}

\subsection{First adiabatic invariant}

\label{sec2-9-1}

Poincaré invariants are generally of little interest in practice unless the curve $\mathcal{C}$ corresponds closely to the actual trajectories of particles. As in previous studies, we consider the case where the curve $\mathcal{C}$ represents a gyration period in phase space. The Poincaré invariant is thus approximated by the associated adiabatic invariant:

\begin{equation}
    \mathcal{I}\simeq\mathcal{J}=\oint_\mathcal{C}\boldsymbol{\mathcal{P}}\cdot\frac{\partial\mathbf{q}}{\partial\varphi}\,\mathrm{d}\varphi
\end{equation}

where $\varphi$ corresponds indeed to the angular parameter of the position $\mathbf{r}_\text{L}(t)$, i.e. $\varphi=\Omega_ct+\varphi_0$. The canonical momentum of a particle of mass $m$ of charge $q$ appears in the following form:

\begin{equation}
    \boldsymbol{\mathcal{P}}=m\mathbf{v}+q\mathbf{A}
\end{equation}

where $\mathbf{A}$ designates the vector potential which satisfies $\nabla\times\mathbf{A}=\mathbf{B}$. The expression of $\mathbf{A}$ in the vicinity of the guiding centre can be written by a Taylor expansion:

\begin{equation}
    \mathbf{A}(\mathbf{r})=\mathbf{A}_0+(\mathbf{r}\cdot\nabla)\mathbf{A}+\mathcal{O}(\mathbf{r}^2)
\end{equation}

It is clear that the next step is to account for a slightly variable magnetic field. In other words, during a gyration period, the magnitude $\|\mathbf{B}\|$ remains approximately constant. The angular parameter $\varphi$ satisfies the following relation:

\begin{equation}
    \mathbf{v}_\text{L}=\Omega_c\diff{\mathbf{r}_\text{L}}{\varphi}\implies\mathrm{d}\boldsymbol{\ell}=\frac{\mathbf{v}_\text{L}}{\Omega_c}\,\mathrm{d}\varphi
\end{equation}

where $\mathrm{d}\boldsymbol{\ell}$ represents a length element belonging to the curve $\mathcal{C}\equiv[0,2\pi)$:

\begin{equation}
    \mathcal{J}=\oint_\mathcal{C}\bigg[m(\mathbf{v}_\mathrm{d}+\mathbf{v}_\text{L})+q(\mathbf{A}_0+(\mathbf{r}_\text{L}\cdot\nabla)\mathbf{A})\bigg]\cdot\frac{m\mathbf{v}_\text{L}}{q\|\mathbf{B}\|}\,\mathrm{d}\varphi
\end{equation}

It is easy to see that the last expression is equivalent to evaluating the mean quantities:

\begin{equation}
    \mathcal{J}=\oint_\mathcal{C}\left[\frac{m^2\mathbf{v}_\mathrm{d}\cdot\mathbf{v}_\text{L}}{q\|\mathbf{B}\|}+\frac{m^2\mathbf{v}_\text{L}^2}{q\|\mathbf{B}\|}+\frac{m\mathbf{A_0}\cdot\mathbf{v}_\text{L}}{\|\mathbf{B}\|}+\frac{m\mathbf{v}_\text{L}}{\|\mathbf{B}\|}\cdot(\mathbf{r_\text{L}}\cdot\nabla)\mathbf{A}\right]\mathrm{d}\varphi
\end{equation}

The first and third terms are zero by orthogonality. The second term is assumed constant, so it comes out of the integral without fail. It remains for us to evaluate the last term:

\begin{equation}
    \mathcal{J}=\frac{2\pi m\mathbf{v}_\text{L}^2}{q\|\mathbf{B}\|}+\frac{2\pi m}{\|\mathbf{B}\|}\left\langle\mathbf{v}_\text{L}\cdot(\mathbf{r}_\text{L}\cdot\nabla)\mathbf{A}\right\rangle
\end{equation}

\begin{equation}
    \langle\mathbf{v}_\text{L}\cdot(\mathbf{r}_\text{L}\cdot\nabla)\mathbf{A}\rangle\equiv\langle v^jx^i\partial_iA_j\rangle\simeq\langle v^ix^j\rangle\partial_iA_j=-\frac{m\mathbf{v}_\text{L}^2}{2q\|\mathbf{B}\|}\left(\nabla\times\mathbf{A}\right)\cdot\mathbf{\hat{B}}=-\frac{m\mathbf{v}_\text{L}^2}{2q}
\end{equation}

\begin{equation}
    \mathcal{J}=\frac{2\pi m^2\mathbf{v}_\text{L}^2}{q\|\mathbf{B}\|}-\frac{2\pi m}{\|\mathbf{B}\|}\frac{m\mathbf{v}_\text{L}^2}{2q}=\frac{\pi m^2\mathbf{v}_\text{L}^2}{q\|\mathbf{B}\|}=2\pi\times\frac{m}{q}\times\frac{m\mathbf{v}_\text{L}^2}{2\|\mathbf{B}\|}
\end{equation}

The reason why we separated the term $m\mathbf{v}_\text{L}^2/2\|\mathbf{B}\|$ is that it represents precisely the magnetic moment of the particle considered. The magnetic moment of any current localization on $\Omega\subset\mathbb{R}^3$ is defined by the following integral. We apply it directly to the particle:

\begin{equation}
    \boldsymbol{\mu}=\frac{1}{2}\iiint_\Omega\mathbf{r}\times\mathbf{J}\,\mathrm{d}V=\frac{1}{2}\mathbf{r}_\text{L}\times q\mathbf{v}_\text{L}=-\frac{q}{2}\frac{m\mathbf{v}_\text{L}^2}{q\|\mathbf{B}\|}\,\mathbf{\hat{B}}=-\frac{m\mathbf{v}_\text{L}^2}{2\mathbf{B}^2}\mathbf{B}
\end{equation}

It then follows that, to the lowest order, the adiabatic invariant is proportional to the magnetic moment $\|\boldsymbol{\mu}\|$. The magnetic moment of a charged particle in a magnetic field is commonly designated by the term \textit{first adiabatic invariant}.

\subsection{Second adiabatic invariant}

\label{sec2-9-2}

The second adiabatic invariant concerns the motion of particles in a magnetic field, as examined in subsection (\ref{sec2-5-2}), subject to the imposed assumptions. We consider a particle of mass $m$ and charge $q$ whose velocity component parallel to the field is $\mathbf{v}_\parallel=(\mathbf{v}\cdot\mathbf{\hat{B}})\,\mathbf{\hat{B}}$. In addition, we will consider the case in which the magnetic field lines are tightly compressed at a point, for example, by a Helmholtz coil. Certain particles, and we will determine which ones, will bounce back in the opposite direction, resulting in a round trip. The bouncing is precisely due to result (\ref{eq2-42}) of subsection (\ref{sec2-5-2}):

\begin{equation}
    \langle\mathbf{F}\cdot\mathbf{\hat{z}}\rangle=-\frac{m\mathbf{v}_\text{L}^2}{2\|\mathbf{B}\|}\frac{\partial B_\text{z}}{\partial z}=-\|\boldsymbol{\mu}\|\frac{\partial B_\text{z}}{\partial z}=\left(\boldsymbol{\mu}\cdot\nabla\right)\mathbf{B}\cdot\mathbf{\hat{z}}
\end{equation}

Conservation of energy and the first adiabatic invariant allow us to write:

\begin{equation}
    \diff{}{t}\left[K=\frac{1}{2}m\left(\mathbf{v}_{\perp}^2+\mathbf{v}_\parallel^2\right)\right]=0\quad\text{et}\quad\diff{\boldsymbol{\mu}}{t}=\diff{}{t}\left(\frac{m\mathbf{v}_\perp^2}{2\mathbf{B}^2}\mathbf{B}\right)=\mathbf{0}
\end{equation}

Consider a confined particle. The component of its velocity parallel to the magnetic field vanishes at the squeezing location. Let $v_\perp$ and $v_\parallel$ be the velocity components when the particle is in the middle between the two boundaries, and let $u_\perp$ be the component perpendicular to the magnetic field $\mathbf{B}_\text{m}$ when it reaches the turnaround point of its trajectory:

\begin{equation}
    \frac{1}{2}m(v_{\parallel}^2+v_\perp^2)\leq\frac{1}{2}mu_\perp^2\quad\text{et}\quad\frac{mv_\perp^2}{2\|\mathbf{B}_0\|}=\frac{mu_\perp^2}{2\|\mathbf{B}_\text{m}\|}\implies\frac{\|\mathbf{B}_0\|}{\|\mathbf{B}_\text{m}\|}\geq\frac{v_\perp^2}{v_\perp^2+v_\parallel^2}
\end{equation}

It is appropriate to introduce the angle $\theta$ between the initial velocity $\mathbf{v}_0$ and the magnetic field $\mathbf{B}_0$, then:

\begin{equation}
    \label{eq2-79}
    \cos(\theta)=\mathbf{\hat{v}}_0\cdot\mathbf{\hat{B}_0}\implies\tan(\theta)=\frac{v_\perp}{v_\parallel}\implies\frac{\|\mathbf{B}_0\|}{\|\mathbf{B}_\text{m}\|}\geq\sin^2(\theta)\implies\theta_c=\sqrt{\arcsin\left(\frac{\|\mathbf{B}_0\|}{\|\mathbf{B}_\text{m}\|}\right)}
\end{equation}

It is easy to see that the condition under which a particle is reflected by the magnetic field depends solely on the angle $\theta$, that is to say, on the ratio between the components perpendicular and parallel to the magnetic field. The critical value is then given by the equality in the last inequality (\ref{eq2-79}). It is appropriate to represent these components using a \textit{loss cone}. It is so named because, if we place the velocity vector of a particle and it is contained within the cone, the particle will not be reflected:

\begin{figure}[H]
    \centering
    \scalebox{1}{\makeatletter
\tikzset{use path/.code=\tikz@addmode{\pgfsyssoftpath@setcurrentpath#1}}
\makeatother

\pgfdeclarelayer{bg}
\pgfsetlayers{bg, main}

\begin{tikzpicture}
    \coordinate (O) at (0, 0, 0);
    
    \draw[thick, black, -Stealth] (0, -3) -- (0, 3) node[anchor=north west] {$v_\perp$};
    \draw[thick, black, -Stealth] (1, 2) -- (-1, -2) node[anchor=north west] {$v_\perp$};
    
    \draw[thick, black, fill=blue!10, dashed] (-4, 0) ellipse (0.5cm and 2cm);
    \draw[thick, black, fill=blue!20, dashed] (4, 0) ellipse (0.5cm and 2cm);

    \draw[thick, black, -Stealth] (-5, 0) -- (5, 0) node[anchor=north] {$v_\parallel$};
    \draw[thick, -Stealth] (2, 0) arc (0:26.75:2);
    \node at (1.33, 0.2875) {$\theta_c$};

    \draw[thick, black, name path=silaz] (-4.5, 2.25) -- (4.5, -2.25);
    \draw[thick, black, name path=uzlaz] (-4.5, -2.25) -- (4.5, 2.25);
    
    \path[name path=L] (-4, 2) -- (-4, -2);
    \path[name path=R] (4, 2) -- (4, -2);
    \path[name intersections={of=L and silaz, by=Lsilaz}];
    \path[name intersections={of=L and uzlaz, by=Luzlaz}];
    \path[name intersections={of=R and uzlaz, by=Ruzlaz}];
    \path[name intersections={of=R and silaz, by=Rsilaz}];
    \path[name intersections={of=silaz and uzlaz, by=mix}];

    \begin{pgfonlayer}{bg}
    \path[intersection segments={of=L and silaz, sequence={L2[reverse]--R2}}, name path=Lsilaz];
    \path[fill=blue!10, intersection segments={of=Lsilaz and uzlaz, sequence={L2--R2}}];
    \path[intersection segments={of=uzlaz and silaz, sequence={L2[reverse]--R2}}, name path=mix];
    \path[fill=blue!10, intersection segments={of=mix and R, sequence={L2--R2}}];
    \end{pgfonlayer}
    
\end{tikzpicture}}
    \caption{\centering Loss cone}
    \label{fig2-1}
\end{figure}

We will next consider a magnetic field \textit{squeezed} at two locations, by two currents, such that the field lines form a solid of revolution. Then, there exists an axial symmetry of the field lines. Particles satisfying (\ref{eq2-79}) will be trapped between the two boundaries. Let $\mathcal{C}$ be the curve tracing the trajectory of the guiding centre. We will tolerate plausible deformations of the curve $\mathcal{C}$. Let $\mathbf{u}$ be the velocity of any point on the curve $\mathcal{C}$. We will impose the following assumption:

\begin{enumerate}
    \item The curve $\mathcal{C}$ undergoes \textit{slow} and \textit{weak} deformations, or none at all, such that the field $\mathbf{u}$ satisfies $\|\mathbf{u}\|\ll r_\text{L}/\tau$, where $r_\text{L}$ is the Larmor radius of a trapped particle and $\tau$ represents the time necessary for the particle to make a complete round trip. This allows us to impose a following bound, regarding the divergence of the field $\mathbf{u}$, which we will need at the end:
    
    \begin{equation}
        \nabla\cdot\mathbf{u}\simeq\|\mathbf{u}\|/\mathbf{r}_L\ll1/\tau
    \end{equation}
\end{enumerate}

The second invariant will be represented by an integral whose derivative with respect to time we will prove to be \textit{approximately} zero. Let $\mathcal{S}$ be a surface where $\mathcal{C}=\partial\mathcal{S}$. Let us immediately apply Stokes' theorem:

\begin{equation}
    \label{eq2-80}
    \mathcal{J}=\oint_\mathcal{C}\mathbf{p}\cdot\mathrm{d}\boldsymbol{\ell}=\iint_\mathcal{S}\left(\nabla\times\mathbf{p}\right)\cdot\mathrm{d}\mathbf{S}\quad\text{où}\quad\mathrm{d}\boldsymbol{\ell}\parallel\mathbf{B}
\end{equation}

We evaluate the time derivative of expression (\ref{eq2-80}). Note that $\mathbf{F}=\partial_t\mathbf{p}$:

\begin{equation}
    \diff{\mathcal{J}}{t}=\diff{}{t}\iint_\mathcal{S}\left(\nabla\times\mathbf{p}\right)\cdot\mathrm{d}\mathbf{S}=\iint_\mathcal{S}\left[\nabla\times\mathbf{F}+\mathbf{u}\left(\nabla\cdot(\nabla\times\mathbf{p})\right)+\nabla\times\left((\nabla\times\mathbf{p})\times\mathbf{u}\right)\right]\cdot\mathrm{d}\mathbf{S}
\end{equation}

Note that the divergence of a curl is always zero, that is, $\nabla\cdot(\nabla\times\mathbf{p})=0$. Since the field lines form a solid of revolution, the surface $\mathcal{S}$ is perpendicular to the magnetic field, so the integration of the first term leads to its vanishing:

\begin{equation}
    \nabla\times(q\mathbf{v}\times\mathbf{B})= q(\underbrace{\nabla\cdot\mathbf{B}}_0)\mathbf{v}+q\mathbf{B}(\nabla\cdot\mathbf{v})=q(\nabla\cdot\mathbf{v})\mathbf{B}\perp\mathrm{d}\mathbf{S}
\end{equation}

\begin{equation}
    \nabla\times\left[(\nabla\times\mathbf{p})\times\mathbf{u}\right] =\nabla\times\left[(\nabla\cdot\mathbf{u})\mathbf{p}-\nabla(\mathbf{p}\cdot\mathbf{u})\right]=\nabla\times(\nabla\cdot\mathbf{u})\mathbf{p}
\end{equation}

\begin{equation}
    \diff{\mathcal{J}}{t}=\iint_\mathcal{S}q(\nabla\cdot\mathbf{v})\underbrace{\mathbf{B}\cdot\mathrm{d}\mathbf{S}}_0+\oint_\mathcal{C}(\nabla\cdot\mathbf{u})\mathbf{p}\cdot\mathrm{d}\boldsymbol{\ell}
\end{equation}

We apply the previously stated hypothesis. The adiabatic invariant is thus quasi-invariant unless the curve $\mathcal{C}$ undergoes uniform deformations, in which case $\nabla\cdot\mathbf{u}=0$, thus maintaining the constancy of the invariant. The second adiabatic invariant is given by the following:

\begin{equation}
    \diff{\mathcal{J}}{t}=\diff{}{t}\oint_\mathcal{C}\mathbf{p}\cdot\mathrm{d}\boldsymbol{\ell}=\oint_\mathcal{C}(\nabla\cdot\mathbf{u})\mathbf{p}\cdot\mathrm{d}\mathbf{S}\simeq0
\end{equation}

\subsection{Third adiabatic invariant}

The third adiabatic invariant concerns the temporal variation of the magnetic field, which was omitted in the previous discussions, i.e., in the analyses of the first and second adiabatic invariants. In this subsection, we will explicitly examine a magnetic field subjected to a temporal variation under certain limiting conditions.

\quad Consider a configuration of a trapping magnetic field. It is appropriate to examine the static effects separately, when the field does not vary, and those that the temporal variation of the field exerts on the dynamics of particles. When we assume that the field does not vary temporally, we refer to this configuration as \textit{static conditions}. A precise description of static conditions is given by the set of following assumptions:

\begin{enumerate}
    \item\label{hyp2-9-3-1} Let $\tau$ be the time necessary for a particle to traverse the curve $\mathcal{C}$ in the preceding considerations. It is more precisely given by the following expression:

    \begin{equation}
        \tau=\oint_\mathcal{C}\frac{\mathrm{d}\ell}{\|\mathbf{v}_\mathrm{d}\|}
    \end{equation}

    Once temporal variation is introduced, the magnetic field subjected to an \textit{adiabatic} variation satisfies the following condition:

    \begin{equation}
        \tau\ll\inf\|\mathbf{B}\|\bigg/\sup\bigg\|\frac{\partial\mathbf{B}}{\partial t}\bigg\|
    \end{equation}

    where the infimum and supremum are taken over the surface $\mathcal{S}$ bounded by the particle's trajectory, in other words by the curve $\mathcal{C}=\partial\mathcal{S}$.
    \item\label{hyp2-9-3-2} The uniformity of the magnetic field in space is very stable. This indeed presents assumption ($\ref{hyp2-5-1-2}$) in subsection (\ref{sec2-5-1}).
\end{enumerate}

Consider the magnetic flux through the surface $\mathcal{S}$ whose boundary is the curve $\mathcal{C}$:

\begin{equation}
    \Phi_\mathbf{B}=\iint_\mathcal{S}\mathbf{B}\cdot\mathrm{d}\mathbf{S}=\oint_\mathcal{C}\mathbf{A}\cdot\mathrm{d}\boldsymbol{\ell}
\end{equation}

\begin{equation}
    \diff{\Phi_\mathbf{B}}{t}=\diff{}{t}\oint_\mathcal{C}\mathbf{A}\cdot\mathrm{d}\boldsymbol{\ell}=\oint_\mathcal{C}\left[\frac{\partial\mathbf{A}}{\partial t}+(\nabla\cdot\mathbf{u})\mathbf{A}\right]\cdot\mathrm{d}\boldsymbol{\ell}
\end{equation}

where $\mathbf{u}$ represents the same velocity field regarding possible deformations of the curve $\mathcal{C}$ as in the previous subsection. This velocity is precisely determined by the variation of the magnetic field over time. Note the Coulomb gauge $\nabla\cdot\mathbf{A}=0$:

\begin{equation}
    \nabla\cdot\mathbf{u}=\nabla\cdot\left(\frac{\mathbf{E}_\text{ind}\times\mathbf{B}}{\mathbf{B}^2}\right)=\nabla\cdot\left(-\frac{\partial\mathbf{A}}{\partial t}\times\frac{\mathbf{B}}{\mathbf{B}^2}\right)=\frac{\partial}{\partial t}\left[\nabla\cdot\left(-\frac{\nabla\mathbf{A}^2}{\mathbf{B}^2}\right)\right]
\end{equation}

Assumption (\ref{hyp2-5-1-2}) of section (\ref{sec2-5-1}) allows us to conclude that this effect of modification of the curve $\mathcal{C}$ is largely negligible compared to the purely temporal variation of the field:

\begin{equation}
    \diff{\Phi_\mathbf{B}}{t}\simeq\oint_\mathcal{C}\frac{\partial\mathbf{A}}{\partial t}\cdot\mathrm{d}\boldsymbol{\ell}=\iint_\mathcal{S}\frac{\partial\mathbf{B}}{\partial t}\cdot\mathrm{d}\mathbf{S}\leq\sup_{\mathbf{B}|_\mathcal{S}}\frac{\partial\|\mathbf{B}\|}{\partial t}S\ll\frac{\inf\|\mathbf{B}\|}{\tau}S\leq\frac{\Phi_\mathbf{B}}{\tau}
\end{equation}

This precisely defines the third adiabatic invariant as the magnetic flux enclosed by the particle's trajectory. It should be noted that, in the purely static case, the magnetic flux is effectively constant.

\section{Violation of adiabatic invariants}

Adiabatic invariants are not entirely conserved quantities, as they are approximations of Poincaré invariants. However, their stability is carefully verified under the aforementioned assumptions when considering a single particle. It is essential to note that the stability of adiabatic invariants in a plasma is considerably lower than that observed for a single particle. Indeed, collisional effects lead to violations of adiabatic invariants. It should be emphasized that the low collisionality of plasmas constitutes a significant advantage, ensuring that our analyses are not invalidated by uncertainty related to collisional effects. Indeed, an interparticle collision significantly modifies the velocities of the two particles involved. Consequently, during a collision, no adiabatic invariant is approximately conserved and undergoes a modification of order of magnitude that is entirely dependent on the collisional conditions. Estimates of these violations, with respect to the first and third adiabatic invariants, cannot be obtained without loss of generality. However, we will provide illustrative examples of various violations of adiabatic invariants.

\subsection{Violation of the second adiabatic invariant}

Consider a collection of charged particles trapped in a magnetic field configuration like that examined in subsection (\ref{sec2-9-2}). Particles belonging to the loss cone illustrated in Figure (\ref{fig2-1}) will escape confinement. This particle escape will cause non-uniformity in particle density in phase space. It will be shown later that this configuration in phase space is unstable, as residual particles will collide. Indeed, collisions serve as carriers of particle diffusion in phase space, thereby leveling and balancing density in this space. This dynamics leads to the entry of particles initially located outside the loss cone, allowing their escape from confinement.

\subsection{Non-adiabatic motion in symmetric geometry}

As we have demonstrated in previous sections, motion characterized by an adiabatic invariant occurs when variations, whether temporal or spatial, in electromagnetic fields during gyration orbits are sufficiently gradual to be treated as analytic. As mentioned previously, when non-adiabatic motion occurs, we do not have the necessary knowledge to develop dynamic models that effectively describe violations of adiabatic invariants. This lack of adequate tools is due to the non-analytic nature of electromagnetic fields in these conditions. Nevertheless, certain situations can be the subject of intuition-rich analysis in non-adiabatic motion. By way of illustration, we will consider a situation where the electromagnetic field is perfectly symmetric with respect to the coordinate $q_i$. It is precisely symmetry that will allow us to develop analytical descriptions of non-adiabatic motion. Symmetry along coordinate $q_i$ ensures that the generalized momentum, denoted $\mathcal{P}_i$, is an exact constant of motion, by virtue of Noether's theorem. In other words, symmetry determines analyticity. The constancy of a motion variable allows the production of a partial, or even complete, solution to the considered non-adiabatic motion problem. Solutions to symmetric problems will help develop intuition that proves valuable when transitioning to non-analytical asymmetric problems.

\subsection{Digression: Stream functions}

When a two-dimensional flow satisfies the continuity equation, $\nabla\cdot\mathbf{u}=0$, it is appropriate to introduce a scalar field, called the \textit{stream function}. A zero-divergence field can always be represented by the curl of a vector potential $\mathbf{A}$ whose stream function is part of the coordinates. Depending on the geometry of the problem, it is possible to distinguish two types of stream functions: Lagrange functions and Stokes functions. First, we study Lagrange stream functions.

\subsubsection{Lagrange Stream Function}

We consider here an irrotational, purely two-dimensional flow in the $xy$ plane. Let $\mathcal{R}^\mathbf{\hat{x}}_{\pi/2}$ be the matrix corresponding to a rotation of $90^\circ$ around the $\mathbf{\hat{x}}$ axis. It is undoubted that a two-dimensional flow $\mathbf{u}$, whether compressible or not, satisfies the following equation:

\begin{equation}
   (\nabla\cdot\mathbf{u})\,\mathbf{\hat{z}}=-\nabla\times(\mathcal{R}^\mathbf{\hat{x}}_{\pi/2}\mathbf{u})
\end{equation}

In case of incompressibility, it follows that the curl of the field $\mathcal{R}^\mathbf{\hat{x}}_{\pi/2}\mathbf{u}$ is exactly zero. By virtue of Stokes' theorem, the integral of the field $\mathcal{R}^\mathbf{\hat{x}}_{\pi/2}\mathbf{u}$ vanishes on any closed curve:

\begin{equation}
    \nabla\times\left(\mathcal{R}^\mathbf{\hat{x}}_{\pi/2}\mathbf{u}\right)=\mathbf{0}\implies\oint_{\partial\mathcal{C}}(\mathcal{R}^\mathbf{\hat{x}}_{\pi/2}\mathbf{u})\cdot\mathrm{d}\boldsymbol{\ell}=0
\end{equation}

The integral being path-independent, it follows, by virtue of the converse of the gradient theorem, that there exists a scalar field satisfying:

\begin{equation}
    \label{eq2-95}
    \mathcal{R}^\mathbf{\hat{x}}_{\pi/2}\mathbf{u}=\nabla\psi
\end{equation}

Having demonstrated the existence of a stream function for an incompressible flow, we will write equation (\ref{eq2-95}) in the following form:

\begin{equation}
    \label{eq2-96}
    \mathbf{u}=\nabla\times\mathbf{A}=\nabla\times\left(\psi\,\mathbf{\hat{z}}\right)=\nabla\times\left(\psi\,\nabla z\right)=\psi \,\underbrace{\nabla\times\mathbf{\hat{z}}}_\mathbf{0}+\nabla\psi\times\mathbf{\hat{z}}=\nabla\psi\times\nabla z
\end{equation}

It follows that the flux through any curve is given by the following:

\begin{equation}
    \label{eq2-98}
    \Phi_\Gamma=\int_\Gamma\mathbf{u}\cdot\mathrm{d}\boldsymbol{\ell}=\int_\Gamma\left(\nabla\psi\times\nabla z\right)\cdot\mathrm{d}\boldsymbol{\ell}=\Delta\psi
\end{equation}

Furthermore, the velocity field and the gradient of the stream function are orthogonal. This observation allows defining a streamline implicitly:

\begin{equation}
    \label{eq2-99}
    \mathbf{u}\cdot\nabla\psi=0\implies\psi(x,y,t)=\text{const. along a streamline}
\end{equation}

\subsubsection{Stokes Stream Function}

\label{sec2-10-3-2}

An axially symmetric flow in cylindrical geometry is certainly not two-dimensional. However, the following symmetry allows us to define another form of stream function:

\begin{equation}
    \label{eq2-100}
    \frac{\partial\mathbf{u}}{\partial\varphi}=\mathbf{0}\implies\mathbf{u}\cdot\boldsymbol{\hat{\varphi}}=\text{const.}=0\,\,\text{according to our needs}
\end{equation}

Like definition (\ref{eq2-99}), the implicit equation of a streamline is given by:

\begin{equation}
    \psi(r,z,t)=\text{const. along a streamline}
\end{equation}

Since the flow is symmetric under rotation by an angle $\varphi$, we determine that the Stokes stream function satisfies the following expression. Nota bene: $\nabla\varphi=\boldsymbol{\hat{\varphi}}/r$:

\begin{equation}
    \label{eq2-101}
    \mathbf{u}=\frac{1}{2\pi}\nabla\psi\times\nabla\varphi
\end{equation}

A flow described by expressions (\ref{eq2-100}) and (\ref{eq2-101}) is called \textit{poloidal}. Like equation (\ref{eq2-98}), the flow flux through a circle of radius $r$ centred at $z$ in the radial plane is given by the following:

\begin{equation}
    \Phi_{\mathcal{C}(r,z)}(t)=\iint_{\mathcal{C}(r,z)}\mathbf{u}\cdot\mathrm{d}\mathbf{S}=\frac{1}{2\pi}\iint_{\mathcal{C}(r,z)}\left(\nabla\psi\times\nabla\varphi\right)\cdot\mathrm{d}\mathbf{S}=\psi\,\biggr|_{(0,z,t)}^{(r,z,t)}
\end{equation}

We isolate the radial component of the field $\mathbf{u}$. Since this field is divergence-free, the radial component at $r=0$ must vanish, so:

\begin{equation}
    u_\text{r}=-\frac{1}{2\pi r}\frac{\partial\psi}{\partial z}=0\implies\frac{\partial\psi}{\partial z}\biggr|_{r=0}=0\implies\psi(0,z,t)=\text{const.}=0\,\,\text{for convenience}
\end{equation}

Thus, the flux through a circle of radius $r$ centred at $z$ in the radial plane is equal to the stream function evaluated on the edge of the circle, in other words, at $(r,z)$:

\begin{equation}
    \Phi_{\mathcal{C}(r,z)}(t)=\psi(r,z,t)
\end{equation}

We will examine two illustrative examples of non-adiabatic motion. These examples concern (i) a sudden temporal reversal and (ii) a sudden spatial reversal of the polarity of an axially symmetric magnetic field with zero azimuthal component. Indeed, the magnetic field being a divergence-free field, it can be described using a Stokes stream function:

\begin{equation}
    \mathbf{B}(r,z,t)=\frac{1}{2\pi}\nabla\psi(r,z,t)\times\nabla\varphi\implies\Phi_{\mathcal{C}(r,z)}=\iint_{\mathcal{C}(r,z)}\mathbf{B}\cdot\mathrm{d}\mathbf{S}=\psi(r,z,t)
\end{equation}

\begin{figure}
    \centering
    \scalebox{1}{\tikzset{
    arw/.style args={#1and#2with step#3}{%
    postaction=decorate,
    decoration={
      markings,
      mark={between positions #1 and #2 step #3 with \arrow{latex}}
    }
  }
}

\begin{tikzpicture}[scale=1.1, z={(-170:10mm)}, x={(-10:10mm)}]

    \tikzset{xzplane/.style={canvas is xz plane at y=#1, very thin}}

    \draw[thick, dashed, -Stealth] (0, -2.5, 0) -- (0, 3.5, 0) node[anchor=east] {$\mathbf{\hat{z}}$};
    
    \foreach \n in {3, ..., 4} {
        \begin{scope}[rotate around y={\n*360/5+4}]
            \coordinate (top) at (0, 2.5, 2.5);
            \coordinate (bottom) at (0, -2.5, 2.5);
            \draw[thick, black, arw=0.1 and 1 with step 50pt] (bottom) .. controls (0, -0.5, 1) and (0, 0.5, 1) .. (top);
            \draw[thick, black, arw=1 and 1 with step 50pt] (top) -- ++(0, 1, 0.8);
        \end{scope}
    }

    \begin{scope}[xzplane=2]
        \draw[thick, black, fill=gray!80, opacity=0.5] (0, 0) circle (2cm);
    \end{scope}

    \draw[thick, black, fill] (0, 2, 0) circle (1.25pt);

    \foreach \n in {5, ..., 7} {
        \begin{scope}[rotate around y={\n*360/5+4}]
            \coordinate (top) at (0, 2.5, 2.5);
            \coordinate (bottom) at (0, -2.5, 2.5);
            \draw[thick, black, arw=0.1 and 1 with step 50pt] (bottom) .. controls (0, -0.5, 1) and (0, 0.5, 1) .. (top);
            \draw[thick, black, arw=1 and 1 with step 50pt] (top) -- ++(0, 1, 0.8);
        \end{scope}
    }

    \begin{scope}[rotate around y={7*360/5-15}]
        \draw[thick, black, fill] (0, 2, 2) circle (1.25pt);
        \draw[thick, black, -Stealth] (0, 2, 0) -- (0, 2, 2) node[midway, above] {$\mathbf{r}$};
        \node at (0, 2, -0.75) {$\psi(r,z,t)$};
        \node at (0, 3, 3) {$\mathbf{B}$};
        \node at (-2, 2, -2) {$\color{gray}\mathcal{C}(r,z)$};
    \end{scope}
    
\end{tikzpicture}}
    \caption{Poloidal $\mathbf{B}$. The flux through $\mathcal{C}(r,z)$ is $\psi(r,z,t)$.}
    \label{fig2-2}
\end{figure}

Like equation (\ref{eq2-96}), the vector potential is given by the following:

\begin{equation}
    \mathbf{A}=\frac{1}{2\pi r}\psi(r,z,t)\,\boldsymbol{\hat{\varphi}}
\end{equation}

Using a Maxwell equation, in the static case, $\nabla\times\mathbf{B}=\mu_0\mathbf{J}$, we determine the electric current giving rise to the considered magnetic field. The current is purely azimuthal, quickly verified:

\begin{equation}
    \mu_0\mathbf{J}\cdot\boldsymbol{\hat{\varphi}}=r\nabla\varphi\cdot(\nabla\times\mathbf{B})=r\nabla\cdot(\mathbf{B}\times\nabla\varphi)=-\frac{r}{2\pi}\nabla\cdot\left(\frac{1}{r^2}\nabla\psi\right)
\end{equation}

\begin{equation}
    \label{eq2-108}
    J_\varphi=-\frac{r}{2\pi\mu_0}\left[\frac{\partial}{\partial r}\left(\frac{1}{r^2}\frac{\partial\psi}{\partial r}\right)+\frac{1}{r^2}\frac{\partial\psi}{\partial z^2}\right]
\end{equation}

The current distribution giving rise to the magnetic field consists, in the most general case, of loops of finite radii $r$ and vertical positions $z$. The axial component of the magnetic field is written:

\begin{equation}
    B_\text{z}=\frac{1}{2\pi r}\frac{\partial\psi}{\partial r}
\end{equation}

The stream function can be expanded in a Taylor series in the vicinity of $r=0$:

\begin{equation}
    \psi(r,z,t)=0+r\frac{\partial\psi}{\partial r}\biggr|_{(0,z,t)}+\frac{1}{2}r^2\frac{\partial^2\psi}{\partial r ^2}\biggr|_{(0,z,t)}+\mathcal{O}(r^3)
\end{equation}

Suppose the term $\partial_r\psi\,(0,z,t)$ is non-zero. Then, in the vicinity of $r=0$, $\psi\sim r$. However, in this case, the right-hand side of equation (\ref{eq2-108}) becomes infinite. This result is obviously unphysical, so we require that the first non-zero term of the Taylor expansion of $\psi$ be the quadratic term. Furthermore, it is evident that the stream function $\psi(r,z,t)$ is positive semi-definite according to our conventions and reaches a maximum. Let $r_\text{max}(z)\equiv r_\text{max}$ such that:

\begin{equation}
    \psi_\text{max}=\sup_{r\in\mathbb{R}_+}\psi(r,z)=\psi(r_\text{max}, z)
\end{equation}

The supremum exists and is reached because the stream function is analytic. The Lagrangian of a particle of mass $m$ and charge $q$ in this cylindrical coordinate system appears in the form:

\begin{equation}
    \mathcal{L}=\frac{1}{2}m\left(\dot{r}^2+r^2\dot{\varphi}^2+\dot{z}^2\right)+qr\dot{\varphi}A_\varphi-q\phi(r,z,t)
\end{equation}

The generalized momentum is then an invariant of motion:

\begin{equation}
    \label{eq2-113}
    \mathcal{P}_\varphi=\frac{\partial\mathcal{L}}{\partial\dot{\varphi}}=mr^2\dot{\varphi}+qrA_\varphi=mr^2\dot{\varphi}+\frac{q}{2\pi}\psi(r,z,t)=\text{const.}
\end{equation}

The Hamiltonian is then given by the following, where we define $\chi(r,z,t)$ as the effective potential:

\begin{equation}
    \mathcal{H}=\frac{1}{2}m\left(\dot{r}^2+r^2\dot{\varphi}^2+\dot{z}^2\right)+\phi(r,z,t)=\frac{1}{2}m(\dot{r}^2+\dot{z}^2)+\frac{(\mathcal{P}_\varphi-q\psi(r,z,t)/2\pi)^2}{2mr^2}+\phi(r,z,t)
\end{equation}

\begin{equation}
    \chi(r,z,t)=\frac{1}{2m}\left[\frac{\mathcal{P}_\varphi-q\psi(r,z,t)/2\pi}{r}\right]^2
\end{equation}

For simplicity, it is appropriate to normalize the effective potential to an arbitrary value. Furthermore, the electrostatic potential will be fixed at zero, due to intuition already developed regarding particle motion in electrostatic fields:

\begin{equation}
    \chi_0=\frac{q\psi_0^2}{8\pi^2mL^2}\implies\frac{\chi(r,z,t)}{\chi_0}=\left[\frac{(2\pi\mathcal{P}_\varphi/q)-\psi(r,z,t)}{(r/L)\psi_0}\right]^2
\end{equation}

The current distribution is assumed static when $t<0$. Since the Lagrangian is time-independent, the energy $\mathcal{H}$ is a constant of motion, associated with $\mathcal{P}_\varphi$ by symmetry.
We consider a particle initially located in the plane $z=z_0$ and radius $r_0<r_\text{max}$. Its motion depends on the sign of $q\psi/\mathcal{P}_\varphi$. We thus distinguish the following two cases:

\begin{enumerate}
    \item $q\psi/\mathcal{P}_\varphi$ is positive. If $2\pi|\mathcal{P}_\varphi|<|q\psi_\text{max}|$, there exist two loci at $r_0<r_\text{max}$ and $r_0'>r_\text{max}$ respectively, where the effective potential $\chi$ is zero:

    \begin{equation}
        \label{eq2-117}
        \mathcal{P}_\varphi=\frac{q}{2\pi}\psi(r_0,z_0)
    \end{equation}

    The effective potential $\chi$ being positive semi-definite, these two loci are also its two global minima, and there exists a local maximum $\chi_\text{max}$ between them. If the energy $\mathcal{H}<\chi_\text{max}$, the particle will be confined in an effective potential well centred around the edge of a surface through which the flux is conditioned by equation (\ref{eq2-117}). The angular velocity $\dot{\varphi}$ changes sign periodically by virtue of equation (\ref{eq2-113}):

    \begin{equation}
        \diff{\varphi}{t}(r,z_0)=\frac{1}{mr^2}\left(\mathcal{P}_\varphi-\frac{q\psi(r,z_0)}{2\pi}\right)
    \end{equation}

    This corresponds precisely to a localized gyration. We will refer to this type of motion by the term \textit{non-encircling} of the $\mathbf{\hat{z}}$ axis.
    
    \item $q\psi/\mathcal{P}_\varphi$ is negative. In this case, $\chi$ reaches zero nowhere, but it reaches a local minimum:

    \begin{equation}
        \left[\left(\mathcal{P}_\varphi-\frac{q}{2\pi}\psi(r,z)\right)\frac{\partial}{\partial r}\left(\frac{\mathcal{P}_\varphi-q\psi(r,z)/2\pi}{r}\right)\right]\Biggr|_{(r_\text{min},z_0)}=0
    \end{equation}

    This equation is satisfied only if the partial derivative is zero:

    \begin{equation}
        \label{eq2-120}
        \frac{\partial}{\partial r}\left(\frac{\mathcal{P}_\varphi-q\psi(r,z)/2\pi}{r}\right)\Biggr|_{(r_\text{min},z_0)}=0\implies\mathcal{P}_\varphi=-\left[\frac{qr^2}{2\pi}\frac{\partial}{\partial r}\left(\frac{\psi(r,z)}{r}\right)\right]\Biggr|_{(r_\text{min},z_0))}
    \end{equation}

    Since the stream function $\psi$ is positive semi-definite, $\psi\sim r^2$ as $r\to0$ and $\psi\to0$ as $r\to\infty$, we determine that $\psi/r\sim r\to0$ as $r\to0$ and $\psi/r\to0$ as $r\to\infty$. Since the function $\psi/r$ is also positive semi-definite, it reaches a maximum in a locus satisfying $r<r_\text{max}$. Let $\sup_{r\leq r_\text{max}}(\psi/r)=\psi(\rho)/\rho$. It should be noted that equation (\ref{eq2-120}) is satisfied only on the condition that $|\mathcal{P}_\varphi|$ is not too large, because the right-hand side of equation (\ref{eq2-120}) reaches a finite maximum. Furthermore, this equation is only valid for points surrounding the maximum of $\psi/r$. Since the function $\psi/r$ increases on the interval $[0,\rho)$, then $\dot{\varphi}$ does not change sign, meaning the particle exerts a purely rotational motion around the $\mathbf{\hat{z}}$ axis. When $r$ is very close to zero, we could approximate equation (\ref{eq2-120}) by the following, clearly similar to equation (\ref{eq2-117}):

    \begin{equation}
        \label{eq2-121}
        \mathcal{P}_\varphi=-\frac{q}{2\pi}\psi(r_\text{min},z_0)
    \end{equation}
    
    Nota bene! This motion is susceptible to undergoing subtle radial variations. Nevertheless, $\dot{\varphi}$ always retains the same sign. We will refer to motions during which $\dot{\varphi}$ retains its sign by the term \textit{encircling} the $\mathbf{\hat{z}}$ axis.
\end{enumerate}

\subsection{Sudden temporal reversal of the magnetic field}

We will now examine a highly non-adiabatic situation in which the current distribution, as considered thus far, undergoes a rapid sign change. In this configuration, the set of fields and fluxes reverses its sign. The energy of the particle cannot be considered a constant of motion, because the Lagrangian explicitly depends on time. However, the symmetry of the problem guarantees the constancy of the azimuthal component of the generalized momentum, $\mathcal{P}_\varphi$. The sign reversal of $\psi$, if performed, can lead to the presence of a minimum of the effective potential $\chi$. A particle not encircling the $\mathbf{\hat{z}}$ axis at a position determined by equation (\ref{eq2-117}) will become encircling. Under the assumption of a particle located near the $\mathbf{\hat{z}}$ axis, by comparison of equations (\ref{eq2-117}) and (\ref{eq2-121}), the radial position of the particle will be conserved. The particle, now encircling the $\mathbf{\hat{z}}$ axis compared to the contrary case before the sign reversal, allows us to observe an increase in its kinetic energy equivalent to the minimum value of $\chi$ of the encircling case.

\quad It would also be appropriate to illustrate the process in question from the perspective of particle drifts. Initially, the particle was frozen by the magnetic field lines, so as to follow the surface characterized by $\psi(r,z)=\text{const.}$ Once the current distribution begins to decrease in intensity, the maximum flux $\psi_\text{max}$ decreases accordingly. The contours of constant $\psi$ inside the circle of radius $r_\text{max}$ move towards this radius where $\psi=\psi_\text{max}$. It is the same for contours of constant $\psi$ outside the aforementioned circle. The annihilation of contours reaching radius $r_\text{max}$ follows. The effect of the $\mathbf{E}\times\mathbf{B}$ drift maintains the motion of particles fixed to a surface of constant flux. Considering a fixed circle of radius $r_0$ at $z=z_0$, the Maxwell-Faraday equation is written as follows:

\begin{equation}
    \oint_{\partial\mathcal{C}(r_0,z_0)}\mathbf{E}\cdot\mathrm{d}\boldsymbol{\ell}=2\pi r_0E_\varphi=-\diff{\Phi_\mathbf{B}}{t}\biggr|_{\mathcal{C}(r_0,z_0)}=-\diff{}{t}\left[\psi(r_0,z_0,t)\right]=-\frac{\partial\psi}{\partial t}\biggr|_{(r_0,z_0)}
\end{equation}

The radial and axial components of the drift velocity give us:

\begin{equation}
    E_\varphi+v_\text{z}B_\text{r}-v_\text{r}B_\text{z}=0\quad\text{with}\quad B_\text{r}=-\frac{1}{2\pi r}\frac{\partial\psi}{\partial z}\quad\text{and}\quad B_\text{z}=\frac{1}{2\pi r}\frac{\partial\psi}{\partial r}
\end{equation}

\begin{equation}
    \frac{\partial\psi}{\partial t}+v_\text{r}\frac{\partial\psi}{\partial r}+v_\text{z}\frac{\partial\psi}{\partial z}\equiv\frac{\partial\psi}{\partial t}+(\mathbf{v}\cdot\nabla)\psi=\frac{\mathrm{d}\psi}{\mathrm{d}t}=0
\end{equation}

The convective derivative being zero, the particle's motion contours a surface of constant $\psi$. It should be noted that the concept of $\mathbf{E}\times\mathbf{B}$ drift collapses completely at the precise moment when $\mathbf{B}=\mathbf{0}$. This presents a failure of the adiabatic approximation. It should be noted that the particle will encircle the $\mathbf{\hat{z}}$ axis following the magnetic field reversal, provided that the stream function undergoes a sign reversal before the particle reaches the surface corresponding to the maximum of the stream function $\psi_\text{max}$. The kinetic energy of the particle increases at the instant when the azimuthal electric field exerts work, when $\mathbf{B}=\mathbf{0}$. Once the sign reversal of $\psi$ is established, the adiabatic approximation is valid again. The polarity of $\psi$ being reversed, the increase in the magnitude of $\psi$ is analogous to an increasing potential well. The particle thus moves away from the surface corresponding to $\psi_\text{min}=-\psi_\text{max}$. Once the reversal is completed, the particles will encircle the $\mathbf{\hat{z}}$ axis and acquire additional kinetic energy, acquired during the instant of sign reversal.

\subsection{Sudden spatial reversal of the magnetic field}

\label{sec2-10-2-2}

We consider an arrangement of two solenoids carrying constant currents, with their magnetic fields opposing each other. Since the currents are constant, the Lagrangian of a particle located inside one of the solenoids does not depend explicitly on time. The energy of the particle is thus a constant of its motion. The given geometry allows us to conclude that the stream function is antisymmetric in $z$, where $z=0$ corresponds to the plane of symmetry between the two solenoids, as illustrated by Figure \ref{fig2-3}.

\quad Consider a particle launched at initial velocity $\mathbf{v}=v_0\,\mathbf{\hat{z}}$ at $z=-L$ and radius $r=a$. Since the velocity component is perpendicular to the $\mathbf{\hat{z}}$ axis, the particle will simply follow a magnetic field line. Once the particle reaches the cusp region, the magnetic field lines begin to curve, as illustrated in Figure (\ref{fig2-3}). Consequently, the particle develops curvature and drifts along the magnetic field gradient. When the particle approaches the plane $z=0$, the magnetic field tends towards zero, which is followed by the collapse of the concept of drift. Despite the complexity of the particle's trajectory in the cusp region, it is possible to determine the condition for a particle to pass through this cusp. The constants of motion are precisely the energy $\mathcal{H}$ and the azimuthal component of the generalized momentum.

\begin{figure}[H]
    \centering
    \scalebox{1}{\tikzset{
    arw/.style args={#1and#2with step#3}{%
    postaction=decorate,
    decoration={
      markings,
      mark={between positions #1 and #2 step #3 with \arrow{Stealth[bend]}}
    }
  }
}

\begin{tikzpicture}

    \draw[dashed, black] (0, -3) -- (0, 3) node[pos=0.38, fill=white] {$z=0$};
    \node at (-2.2, 2.65) {$\mathbf{B}$};

    \foreach \n in {-6, ..., -2} {
        \draw[thick, black] (\n, 2) circle (0.2cm);
        \draw[thick, black, fill] (\n, 2) circle (1.25pt);
        \draw[thick, black] (\n, -2) circle (0.2cm);
        \draw[thick, black] ({\n+0.2/sqrt(2)}, {-2+0.2/sqrt(2)}) -- ({\n-0.2/sqrt(2)}, {-2-0.2/sqrt(2)});
        \draw[thick, black] ({\n-0.2/sqrt(2)}, {-2+0.2/sqrt(2)}) -- ({\n+0.2/sqrt(2)}, {-2-0.2/sqrt(2)});
    }

    \foreach \n in {2, ..., 6} {
        \draw[thick, black] (\n, 2) circle (0.2cm);
        \draw[thick, black, fill] (\n, -2) circle (1.25pt);
        \draw[thick, black] (\n, -2) circle (0.2cm);
        \draw[thick, black] ({\n+0.2/sqrt(2)}, {2+0.2/sqrt(2)}) -- ({\n-0.2/sqrt(2)}, {2-0.2/sqrt(2)});
        \draw[thick, black] ({\n-0.2/sqrt(2)}, {2+0.2/sqrt(2)}) -- ({\n+0.2/sqrt(2)}, {2-0.2/sqrt(2)});
    }

    \draw[thick, black, arw=0.1 and 1 with step 50pt] (-6.5, 1.3333) .. controls (0, 1.3333) and (-1, 1.5) .. (-2, 3);
    \draw[thick, black, arw=0.25 and 1 with step 50pt] (-6.5, 0.6667) .. controls (0.25, 0.6667) and (-0.25, 1) .. (-0.75, 3);
    \draw[thick, black, arw=0.1 and 1 with step 50pt] (-6.5, 0) -- (0, 0);
    \draw[thick, black, arw=0.25 and 1 with step 50pt] (-6.5, -0.6667) .. controls (0.25, -0.6667) and (-0.25, -1) .. (-0.75, -3);
    \draw[thick, black, arw=0.1 and 1 with step 50pt] (-6.5, -1.3333) .. controls (0, -1.3333) and (-1, -1.5) .. (-2, -3);

    \draw[thick, black, arw=0.1 and 1 with step 50pt] (6.5, 1.3333) .. controls (0, 1.3333) and (1, 1.5) .. (2, 3);
    \draw[thick, black, arw=0.25 and 1 with step 50pt] (6.5, 0.6667) .. controls (-0.25, 0.6667) and (0.25, 1) .. (0.75, 3);
    \draw[thick, black, arw=0.1 and 1 with step 50pt] (6.5, 0) -- (0, 0);
    \draw[thick, black, arw=0.25 and 1 with step 50pt] (6.5, -0.6667) .. controls (-0.25, -0.6667) and (0.25, -1) .. (0.75, -3);
    \draw[thick, black, arw=0.1 and 1 with step 50pt] (6.5, -1.3333) .. controls (0, -1.3333) and (1, -1.5) .. (2, -3);

    \draw[thick, blue, -{Stealth[bend]}] (-6.5, 1) to[out=0, in=135] (-5.575, 0.725);
    \draw[thick, blue, arw=0.4 and 1 with step 50pt] (-5.5, 0.625) to[out=-45, in=225] (-3.5, 0.7) to[out=45, in = 145, looseness=0.8] (-1.4, 1);
    \draw[thick, blue, arw=0.6 and 1 with step 50pt] (-1.28, 0.87) .. controls (-1, 0.5) and (0, 1) .. (-0.5, 1.5);
    \draw[thick, blue, arw=0.4 and 1 with step 50pt] (-0.5, 1.5) .. controls (-2, 3) and (1, 4) .. (1.4, 1);
    \draw[thick, blue, arw=0 and 1 with step 50pt] (1.4, 1) .. controls (2, -5) and (3, 3) .. (3.4, 1);
    \draw[thick, blue, arw=0.3 and 1 with step 70pt] (3.4, 1) .. controls (4.2, -4.5) and (4.75, 2.5) .. (5.1, 0.85);
    \draw[thick, blue, fill] (5.1, 0.85) circle (1pt);
    
\end{tikzpicture}}
    \vspace*{-2cm}
    \caption{Lines of $\mathbf{B}$ and trajectory of a particle crossing the cusp in blue.}
    \label{fig2-3}
\end{figure}

\begin{equation}
    \label{eq2-125}
    \mathcal{P}_\varphi=mr^2\dot{\varphi}+\frac{q}{2\pi}\psi(r,z,t)=\frac{q}{2\pi}\psi(a,-L)=\text{const.}
\end{equation}

\begin{equation}
    \label{eq2-126}
    \mathcal{H}=\frac{\mathbf{p}^2}{2m}=\frac{1}{2}mv_0^2=\text{const.}
\end{equation}

The minimum energy to overcome the cusp corresponds to motion encircling the $\mathbf{\hat{z}}$ axis and zero radial and axial components. Equations (\ref{eq2-125}) and (\ref{eq2-126}) give us the sought condition:

\begin{equation}
    \label{eq2-127}
    v_0^2\geq r^2\dot{\varphi}^2=\left(\frac{q}{2\pi mr}\right)^2\left[\psi(a,-L)-\psi(r,z)\right]^2
\end{equation}

Depending on the value of $v_0$, equation (\ref{eq2-127}) will give us several solutions for $r$. Imaginary and negative solutions are unphysical. This corresponds to the particle's inability to penetrate the cusp. It will return towards the interior of the first solenoid. However, positive real solutions correspond to a genuine passage through the cusp, illustrated in figure (\ref{fig2-3}).
